        %%******************************************%%
        %%                                          %%
        %%        Modello di tesi di laurea         %%
        %%            di Andrea Giraldin            %%
        %%                                          %%
        %%             2 novembre 2012              %%
        %%                                          %%
        %%******************************************%%


% I seguenti commenti speciali impostano:
% 1.
% 2. PDFLaTeX come motore di composizione;
% 3. tesi.tex come documento principale;
% 4. il controllo ortografico italiano per l'editor.

% !TEX encoding = UTF-8
% !TEX TS-program = pdflatex
% !TEX root = tesi.tex
% !TEX spellcheck = it-IT

% PDF/A filecontents
\RequirePackage{filecontents}
\begin{filecontents*}{\jobname.xmpdata}
  \Title{Document’s title}
  \Author{Author’s name}
  \Language{it-IT}
  \Subject{The abstract, or short description.}
  \Keywords{keyword1\sep keyword2\sep keyword3}
\end{filecontents*}

\documentclass[10pt,                    % corpo del font principale
               a4paper,                 % carta A4
               twoside,                 % impagina per fronte-retro
               openright,               % inizio capitoli a destra
               english,
               italian,
               "xcolor=table"
               ]{book}

%**************************************************************
% Importazione package
%**************************************************************

\PassOptionsToPackage{dvipsnames}{xcolor} % colori PDF/A

\usepackage{colorprofiles}

\usepackage[a-2b,mathxmp]{pdfx}[2018/12/22]
                                        % configurazione PDF/A
                                        % validare in https://www.pdf-online.com/osa/validate.aspx

%\usepackage{amsmath,amssymb,amsthm}    % matematica

\usepackage[T1]{fontenc}                % codifica dei font:
                                        % NOTA BENE! richiede una distribuzione *completa* di LaTeX

\usepackage[utf8]{inputenc}             % codifica di input; anche [latin1] va bene
                                        % NOTA BENE! va accordata con le preferenze dell'editor

\usepackage[english, italian]{babel}    % per scrivere in italiano e in inglese;
                                        % l'ultima lingua (l'italiano) risulta predefinita

\usepackage{bookmark}                   % segnalibri

\usepackage{caption}                    % didascalie

\usepackage{chngpage,calc}              % centra il frontespizio

\usepackage{csquotes}                   % gestisce automaticamente i caratteri (")

\usepackage{emptypage}                  % pagine vuote senza testatina e piede di pagina

\usepackage{epigraph}			% per epigrafi

\usepackage{eurosym}                    % simbolo dell'euro

%\usepackage{indentfirst}               % rientra il primo paragrafo di ogni sezione

\usepackage{graphicx}                   % immagini

\usepackage{hyperref}                   % collegamenti ipertestuali

\usepackage[binding=5mm]{layaureo}      % margini ottimizzati per l'A4; rilegatura di 5 mm

\usepackage{listings}                   % codici

\usepackage{microtype}                  % microtipografia

\usepackage{mparhack,fixltx2e,relsize}  % finezze tipografiche

\usepackage{nameref}                    % visualizza nome dei riferimenti
\usepackage[font=small]{quoting}        % citazioni

\usepackage{subfig}                     % sottofigure, sottotabelle

\usepackage[italian]{varioref}          % riferimenti completi della pagina

\usepackage{booktabs}                   % tabelle
\usepackage{tabularx}                   % tabelle di larghezza prefissata
\usepackage{longtable}                  % tabelle su più pagine
\usepackage{ltxtable}                   % tabelle su più pagine e adattabili in larghezza
\usepackage{placeins}                   % per \FloatBarrier
\usepackage{mdframed}                   % per contorni immagini
\usepackage[toc, acronym]{glossaries}   % glossario
                                        % per includerlo nel documento bisogna:
                                        % 1. compilare una prima volta tesi.tex;
                                        % 2. eseguire: makeindex -s tesi.ist -t tesi.glg -o tesi.gls tesi.glo
                                        % 3. eseguire: makeindex -s tesi.ist -t tesi.alg -o tesi.acr tesi.acn
                                        % 4. compilare due volte tesi.tex.

\usepackage[backend=biber,style=verbose-ibid,hyperref,backref]{biblatex}
\usepackage{colortbl}
\usepackage{array,booktabs}
\usepackage{listings}
\lstdefinelanguage{JavaScript}{
  %keywords={typeof, new, true, false, catch, function, return, null, catch, switch, var, if, in, while, do, else, case, break},
  %keywordstyle=\color{blue}\bfseries,
  keywords={string, number, Date, typeof, new, true, false, catch, function, return, null, catch, switch, var, if, in, while, do, else, case, break, let},
  keywordstyle=\color{Plum}\bfseries,
  ndkeywords={class, export, boolean, throw, implements, import, this, interface, describe, it},
  ndkeywordstyle=\color{NavyBlue}\bfseries,
  %skewywords={string, numer, Date},
  sensitive=false,
  comment=[l]{//},
  morecomment=[s]{/*}{*/},
  commentstyle=\color{purple}\ttfamily,
  stringstyle=\color{red}\ttfamily,
  morestring=[b]',
  morestring=[b]"
}
\lstset{
   language=JavaScript,
   extendedchars=true,
   basicstyle=\footnotesize\ttfamily,
   showstringspaces=false,
   showspaces=false,
   numbers=left,
   numberstyle=\footnotesize,
   numbersep=9pt,
   tabsize=2,
   breaklines=true,
   showtabs=false,
   captionpos=b
}

                                        % eccellente pacchetto per la bibliografia;
                                        % produce uno stile di citazione autore-anno;
                                        % lo stile "numeric-comp" produce riferimenti numerici
                                        % per includerlo nel documento bisogna:
                                        % 1. compilare una prima volta tesi.tex;
                                        % 2. eseguire: biber tesi
                                        % 3. compilare ancora tesi.tex.

\input{tesi-config}                     % file con le impostazioni personali

\definecolor{mygrey1}{HTML}{BCBCBC}
\definecolor{mygrey2}{HTML}{CFCFCF}
%\newcommand{\centered}[1]{\begin{tabular}{l} #1 \end{tabular}}
\begin{document}
%**************************************************************
% Materiale iniziale
%**************************************************************
\frontmatter
\input{inizio-fine/frontespizio}
\input{inizio-fine/colophon}
% !TEX encoding = UTF-8
% !TEX TS-program = pdflatex
% !TEX root = ../tesi.tex

%**************************************************************
% Dedica
%**************************************************************
\cleardoublepage
\phantomsection
\thispagestyle{empty}
\pdfbookmark{Dedica}{Dedica}

\vspace*{3cm}

\begin{center}
De omnibus dubitandum. \\ \medskip
--- Cartesio
\end{center}

\medskip

% !TEX encoding = UTF-8
% !TEX TS-program = pdflatex
% !TEX root = ../tesi.tex

%**************************************************************
% Sommario
%**************************************************************
\cleardoublepage
\phantomsection
\pdfbookmark{Sommario}{Sommario}
\begingroup
\let\clearpage\relax
\let\cleardoublepage\relax
\let\cleardoublepage\relax

\chapter*{Sommario}
Le moderne Web Application adottano un disaccoppiamento stretto tra client e backend.\\\\
Soluzioni architetturali come REST garantiscono di poter realizzare API di utilizzo generale fruibili da svariati client, dai browser sino alle più moderne applicazioni mobili.\\\\
Pur essendo REST lo standard de-facto per la scrittura di Web API, esso presenta alcune debolezze che nuovi strumenti sorti negli utlimi anni cercano di superare.\\\\
GraphQL è certamente una delle più recenti e popolari tecnologie che il mercato dell'Information Technology ci mette a disposizione per la realizzazione di Web API.\\\\
Obbiettivo di questa tesi è quello di realizzare una comparazione tra le tecnologie di data fetching REST e GraphQL individuando i vantaggi nell'adozione di una tecnologia rispetto all'altra; le caratteristiche dei due approcci emergono chiaramente durante lo sviluppo e migrazione di una applicazione da un approccio all'altro.
\endgroup

\vfill

% !TEX encoding = UTF-8
% !TEX TS-program = pdflatex
% !TEX root = ../tesi.tex

%**************************************************************
% Ringraziamenti
%**************************************************************
\cleardoublepage
\phantomsection
\pdfbookmark{Ringraziamenti}{ringraziamenti}

\begingroup
\let\clearpage\relax
\let\cleardoublepage\relax
\let\cleardoublepage\relax

\chapter*{Ringraziamenti}

\noindent \textit{Sentiti ringraziamenti al Prof. Paolo Baldan che mi ha seguito durante lo stage e la redazione della tesi.}\\

\noindent \textit{Ringrazio anche il tutor aziendale Fabio Pallaro per la disponibilità che mi ha dimostrato durante tutto lo svolgimento dello stage.}\\
\bigskip

\noindent\textit{\myLocation, \myTime}
\hfill \myName

\endgroup

\input{inizio-fine/indici}
\cleardoublepage

%**************************************************************
% Materiale principale
%**************************************************************
\mainmatter
% !TEX encoding = UTF-8
% !TEX TS-program = pdflatex
% !TEX root = ../tesi.tex

%**************************************************************
\chapter{Introduzione}
\label{introduzione}
%**************************************************************
\intro{Il capitolo seguente introduce all'azienda e illustra la struttura del documento.}\\
%**************************************************************
\section{L'azienda}
\FloatBarrier
\begin{figure}[!h]
\centering
\includegraphics[width=0.4\linewidth]{immagini/synclabLogo.jpeg}
%\caption{Logo Synclab.}
\label{controller-service-repository}
\end{figure}
\FloatBarrier
SyncLab S.r.l è un azienda che si occupa di \gls{ict} ed opera in molteplici settori quali: industriale, sanitario, finanziario, logistico, di telecomunicazioni e infine sanitario.\\\\
Nata a Napoli nel 2002 è al giorno d'oggi divisa in sei sedi operative nelle città di Napoli, Roma, Milano, Como, Verona e infine Padova. Conta più di 300 dipendenti e vanta oltre 150 clienti.

\end{figure}

%**************************************************************
\section{Introduzione al progetto e alla soluzione scelta}
L'idea alla base dello stage è quella di effettuare la migrazione di un applicativo aziendale dal protocollo REST a GraphQL. Lo scopo è quello di analizzare la nuova tecnologia mai affrontata dall'azienda e valutare, attraverso la realizzazione di una analisi di comparazione, quale protocollo potesse soddisfare meglio le esigenze dell'applicativo aziendale.\\ \\
Prima di eseguire la migrazione sull'applicativo aziendale è stato scelto di realizzare un prototipo di gestionale aziendale per familiarizzare con le tecnologie e comprendere quale fosse l'approccio corretto per affrontare la migrazione dell'applicativo aziendale al meglio.\\
Trattandosi di un applicativo composto da backend e frontend, il prototipo è stato realizzato seguendo la medesima architettura al fine di comprenderne il funzionamento e fare pratica. É stato implementato inizialmente con REST API per poi effettuarne la migrazione a GraphQL API.\\ \\
L'applicativo aziendale scelto sul quale è stato realizzato il processo di migrazione si chiama SushiLab. L'applicativo SushiLab, pensato e realizzato da alcuni stagisti nei mesi che precedevano il mio stage, mira alla gestione delle ordinazioni nei ristoranti di sushi nella formula \gls{ayce}. Dunque dopo aver analizzato l'applicativo e compreso il suo funzionamento ne è stata effettuata la migrazione.\\ \\
Al termine della migrazione, prendendo atto di quanto compreso durante la progettazione, lo sviluppo e la migrazione del prototipo e di SushiLab, è stata effettuata una analisi comparativa su tutti gli aspetti che differenziano i due protocolli per la realizzazione di API.
\section{Descrizione del prodotto ottenuto}
Al termine dello stage è stato realizzato:
\begin{itemize}
  \item \textbf{Prototipo} composto da backend e frontend è stato realizzato in due versioni:
  \begin{itemize}
    \item \textit{Versione REST}: il prototipo è stato realizzato prima con REST API;
    \item \textit{Versione GraphQL}: a partire dalla versione con REST API è stata effettuata la migrazione e dunque è stato ottenuta una versione funzionante con GraphQL API;
  \end{itemize}
  \item \textbf{SushiLab con GraphQL API}: è stato realizzata la migrazione dell'applicativo aziendale dalla sua versione con REST API, alla versione funzionante con GraphQL API;
  \item \textbf{Analisi comparativa}: è stato realizzato un documento dettagliato sulla comparazione dei due protocolli REST e GraphQL nel quale sono stati approfonditi sia gli aspetti prettamente teorici che quelli più pratici.
\end{itemize}
\newpage
\section{Strumenti e tecnologie utilizzati}
\subsection*{IntelliJ IDEA}
\begin{figure}[ht]
  \begin{minipage}[h]{0.3\linewidth}
    \centering
    \includegraphics[width=0.9\linewidth]{immagini/IntelliJ.png}
  \end{minipage}
  \begin{minipage}[!h]{0.7\linewidth}
    \textit{IntelliJ IDEA} è un \gls{ide} sviluppato da \textit{JetBrains} che agevola nello sviluppo in \gls{jv}.
    É stato utilizzato nello sviluppo del backend del prototipo e durante la migrazione lato backend sia del prototipo che di SushiLab.
  \end{minipage}
\end{figure}
\subsection*{Visual Studio Code}
\begin{figure}[ht]
  \begin{minipage}[h]{0.3\linewidth}
    \centering
    \includegraphics[width=0.4\linewidth]{immagini/visualStudioCode.png}
  \end{minipage}
  \begin{minipage}[!h]{0.7\linewidth}
    \textit{Visual Studio Code} è un \gls{ide} sviluppato da \textit{Microsoft} che fornisce molte feature per agevolare lo sviluppo.
    É stato utilizzato nello sviluppo del frontend del prototipo e durante la migrazione lato frontend sia del prototipo che di SushiLab.
  \end{minipage}
\end{figure}
\subsection*{Atom}
\FloatBarrier
\begin{figure}[!h]
  \begin{minipage}[h]{0.3\linewidth}
    \centering
    \includegraphics[width=0.5\linewidth]{immagini/atom.png}
  \end{minipage}
  \begin{minipage}[!h]{0.7\linewidth}
    \textit{Atom} è un editor di testo open source. É stato utilizzato durante la stesura con linguaggio \gls{ltx} del documento di analisi comparativa e della documentazione del prototipo e di SushiLab.
  \end{minipage}
\end{figure}
\FloatBarrier
\subsection*{Spring Boot Framework}
\FloatBarrier
\begin{figure}[!h]
  \begin{minipage}[h]{0.3\linewidth}
    \centering
    \includegraphics[width=0.4\linewidth]{immagini/spring.png}
  \end{minipage}
  \begin{minipage}[!h]{0.7\linewidth}
    \textit{Spring Boot} è un framework open source per lo sviluppo di applicazioni in \gls{jv} strutturate a microservizi. Associati ad esso troviamo molti altri moduli e librerie ampiamente utilizzate nel corso dello stage poiché se integrati permettono di agevolare lo sviluppo. Utilizzato per lo sviluppo e migrazione del prototipo e di SushiLab.
  \end{minipage}
\end{figure}
\FloatBarrier
\subsection*{Angular Framework}
\FloatBarrier
\begin{figure}[!h]
  \begin{minipage}[h]{0.3\linewidth}
    \centering
    \includegraphics[width=0.5\linewidth]{immagini/angular.png}
  \end{minipage}
  \begin{minipage}[!h]{0.7\linewidth}
    \textit{Angular} è un framework open source per lo sviluppo di applicazioni web. Permette di realizare applicazioni multipiattaforma che si adattano facilmente a diversi dispositivi. Utilizzato per lo sviluppo e la migrazione del frontend nel prototipo e in SushiLab.
  \end{minipage}
\end{figure}
\FloatBarrier
\subsection*{GitHub}
\begin{figure}[ht]
  \begin{minipage}[h]{0.3\linewidth}
    \centering
    \includegraphics[width=0.5\linewidth]{immagini/github.png}
  \end{minipage}
  \begin{minipage}[!h]{0.7\linewidth}
    \textit{GitHub} è un servizio di hosting per progetti software. É stato utilizzato per effettuare i salvataggi su cloud del codice e della documentazione sia per quanto riguarda il prototipo che per SushiLab. Permette di tener traccia di tutte le versioni gestite con lo strumento di versionamento Git.
  \end{minipage}
\end{figure}
\subsection*{JUnit 5}
\FloatBarrier
\begin{figure}[!h]
  \begin{minipage}[h]{0.3\linewidth}
    \centering
    \includegraphics[width=0.6\linewidth]{immagini/JUnit.png}
  \end{minipage}
  \begin{minipage}[!h]{0.7\linewidth}
    \textit{JUnit 5} è un framework di unit testing per il linguaggio di programmazione \gls{jv}. É stato ampiamente utilizzato per effettuare i test di unità nel backend sia del prototipo che dell'applicativo SushiLab.
  \end{minipage}
\end{figure}
\FloatBarrier
\subsection*{Jasmine}
\FloatBarrier
\begin{figure}[!h]
  \begin{minipage}[h]{0.3\linewidth}
    \centering
    \includegraphics[width=0.4\linewidth]{immagini/jasmine.png}
  \end{minipage}
  \begin{minipage}[!h]{0.7\linewidth}
    \textit{JUnit 5} è un framework di test open source. É stato utilizzato in Angular per effettuare i test nel frontend sia del prototipo che dell'applicativo SushiLab.
  \end{minipage}
\end{figure}
\FloatBarrier
\subsection*{Trello}
\FloatBarrier
\begin{figure}[!h]
  \begin{minipage}[h]{0.3\linewidth}
    \centering
    \includegraphics[width=0.35\linewidth]{immagini/trello.png}
  \end{minipage}
  \begin{minipage}[!h]{0.7\linewidth}
    \textit{Trello} è un software gestionale. É stato utilizzato per tenere traccia di tutti gli sviluppi dello stage nel corso delle varie settimane.
  \end{minipage}
\end{figure}
\FloatBarrier
\subsection*{Discord}
\FloatBarrier
\begin{figure}[!h]
  \begin{minipage}[h]{0.3\linewidth}
    \centering
    \includegraphics[width=0.7\linewidth]{immagini/discord.png}
  \end{minipage}
  \begin{minipage}[!h]{0.7\linewidth}
    \textit{Discord} è una piattaforma di messaggistica e distribuzione digitale per la comunicazione. É stata utilizzata per la comunicazione con il tutor aziendale e gli altri stagisti.
  \end{minipage}
\end{figure}
\FloatBarrier
%**************************************************************
\section{Organizzazione del testo}
\begin{description}
    \item Il {\hyperref[descrizione-stage]{secondo capitolo}} descrive l'analisi preventiva dei rischi, obbiettivi e infine la pianificazione del lavoro;

    \item Il {\hyperref[protocolli-trasmissione-dati]{terzo capitolo}} introduce ai protocollo di modellazione e trasferimento, tra questi approfondisce REST e GraphQL;

    \item Il {\hyperref[casi-uso]{quarto capitolo}} mostra la realizzazione del prototipo con relativa migrazione e la migrazione dell'applicativo SushiLab;

    \item Il {\hyperref[analisi-comparativa]{quinto capitolo}} descrive dettagliatamente l'analisi comparativa tra i protocolli REST e GraphQL;

    \item Il {\hyperref[conclusioni]{sesto capitolo}} rirpota le conclusioni finali riguardo lo stage effettuato.
\end{description}

Riguardo la stesura del testo, relativamente al documento sono state adottate le seguenti convenzioni tipografiche:
\begin{itemize}
	\item gli acronimi, le abbreviazioni e i termini ambigui o di uso non comune menzionati vengono definiti nel glossario, situato alla fine del presente documento;
	\item per la prima occorrenza dei termini riportati nel glossario viene utilizzata la seguente nomenclatura: \emph{parola}\glsfirstoccur;
\end{itemize}
             % Introduzione
% !TEX encoding = UTF-8
% !TEX TS-program = pdflatex
% !TEX root = ../tesi.tex

%**************************************************************
\chapter{Descrizione dello stage}
\label{descrizione-stage}
%**************************************************************

\intro{Nel capitolo viene esposta l'organizzazione e la pianificazione dello stage}\\
%**************************************************************
\section{Requisiti e obiettivi}
\label{obiettivi-previsti}
\subsection*{Notazione dei requisiti e degli obbiettivi}
Si farà riferimento ai requisiti secondo le seguenti notazioni:
\begin{itemize}
  \item O per i requisiti obbligatori, vincolanti in quanto obiettivo primario richiesto dal committente;
  \item D per i requisiti desiderabili, non vincolanti o strettamente necessari, ma dal riconoscibile valore aggiunto;
  \item F per i requisiti facoltativi, rappresentanti valore aggiunto non strettamente competitivo.
\end{itemize}
Le sigle indicate saranno seguita da dei numeri sequenziali che identificano i vari requisiti.
\subsection*{Obiettivi fissati}
É stato previsto lo svolgimento dei seguenti obiettivi:
\begin{itemize}
  \item \textbf{Obbligatori}:
  \begin{itemize}
    \item \underline{O01}: realizzazione dello studio con relativa relazione sulle tecnologie di trasferimento e model- lazione dati REST API e GraphQL;
    \item \underline{O02}: studio dei framework Spring e Angular;
    \item \underline{O03}: sviluppo di un prototipo di web application con back-end Spring, front-end Angular e
comunicazione tramite linguaggio GraphQL;
    \item \underline{O04}: realizzazione migrazione da tecnologia REST API a GraphQL di una web application aziendale;
  \end{itemize}
  \item \textbf{Desiderabili}:
  \begin{itemize}
    \item \underline{D01}: realizzazione di una relazione che riporti un’analisi approfondita e le differenze tra le due tecnologie REST API e GraphQL nel caso d’uso specifico della web application scelta per la migrazione;
  \end{itemize}
  \item \textbf{Facoltativi}:
  \begin{itemize}
    \item \underline{F01}: realizzazione della migrazione su ulteriori web application per valutare le tecnologie REST API e GraphQL su ulteriori casi d’uso;
  \end{itemize}
\end{itemize}
%**************************************************************
\section{Analisi preventiva dei rischi}
Nella fase iniziale dello stage sono state individuate alcue possibili criticità che sarebbero potute sorgere durante lo svolgimento dello stage. Per questo motivo sono state individuate alcune soluzioni per evitare tali rischi. Le principali criticità individuate sono:
\begin{itemize}
  \item \textbf{Conoscenza superficiale delle tecnologie utilizzate}: la maggior parte delle tecnologie utilizzate durante lo stage sono tecnologie a me sconosciute. La soluzione individuata prevede che le prime settimane dello stage siano dedicate principalmente allo studio delle seguenti tecnologie:
  \begin{itemize}
    \item \textit{Spring Boot Framework};
    \item \textit{Angular Framework};
  \end{itemize}
  \item \textbf{Comprensione dell'applicativo SushiLab}: non avendo partecipato allo sviluppo di SushiLab potrebbe risultare complesso comprenderne le logiche. La soluzione individuata prevede lo studio dell'applicativo consultando sia la documentazione che il codice sorgente.
\end{itemize}
\section{Pianificazione}
\label{pianificazione}
Inizialmente il carico di lavoro è stato distribuito sulle otto settimane disponibili di stage. Segue la suddivisione:
\begin{itemize}
    \item \textbf{Prima Settimana (40 ore)}
    \begin{itemize}
        \item Incontro con le persone coinvolte nel progetto per discutere i requisiti e le richieste relativamente al sistema da sviluppare;
        \item Verifica credenziali e strumenti di lavoro assegnati;
        \item Ripasso Java Standard Edition e tool di sviluppo (IDE ecc.);
        \item Studio teorico dell’architettura a microservizi: passaggio da monolite ad architetture a microservizi con pro e contro;
        \item Ripasso principi della buona programmazione (SOLID, CleanCode);
        \item Ripasso concetti Web (Servlet, servizi Rest, Json ecc.).
    \end{itemize}
    \item \textbf{Seconda Settimana (40 ore)}
    \begin{itemize}
        \item Studio Spring Core/Spring Boot;
        \item Studio ORM, in particolare il framework Spring Data JPA;
        \item Studio servizi REST e framework Spring Data REST;
        \item Studio linguaggio GraphQL.
    \end{itemize}
    \item \textbf{Terza Settimana (40 ore)}
    \begin{itemize}
        \item Studio linguaggio TypeScript;
        \item Studio Framework Angular.
    \end{itemize}
    \item \textbf{Quarta Settimana (40 ore)}
    \begin{itemize}
        \item Completamento formazione su Angular.
    \end{itemize}
    \item \textbf{Quinta Settimana (40 ore)}
    \begin{itemize}
        \item Realizzazione di un mini prototipo web application con back-end Spring, front-end Angular con comunicazione tramite linguaggio GraphQL.
    \end{itemize}
    \item \textbf{Sesta Settimana (40 ore)}
    \begin{itemize}
        \item Inizio studio e progettazione migrazione progetto didattico REST versus GraphQL;
        \item Sviluppo migrazione.
    \end{itemize}
    \item \textbf{Settima Settimana (40 ore)}
    \begin{itemize}
        \item Proseguimento sviluppo migrazione.
    \end{itemize}
    \item \textbf{Ottava Settimana (40 ore)}
    \begin{itemize}
        \item Fine sviluppi;
        \item Validazione prodotto con stakeholders.
    \end{itemize}
\end{itemize}
\section{Organizzazione dello stage}
\subsection*{Lavoro in sede e Smart Working}
L'organizzazione dello stage è stata accordata con il tutor aziendale. La sede SyncLab di Padova non è aperta tutti i giorni della settimana, per questo motivo sono previsti anche dei giorni di lavoro in modalità Smart Working. Dunque l'organizzazione è stata accordata come segue:
\begin{itemize}
  \item \textbf{Almeno un giorno in sede}: non tutti i giorni della settimana il tutor aziendale è presente in sede. Per questo motivo è stato accordato di lavorare in sede almeno un giorno della settimana in cui il tutor fosse presente. Per sfruttare al massimo l'esperienza in sede ho cercato di frequentare la sede più di una volta a settimana;
  \item \textbf{Lavoro in Smart Working}: nei giorni in cui non si frequentava la sede il lavoro è stato realizzato in modalità Smart Working. In ogni caso la disponibilità del tutor è stata garantita anche in questa modalità, infatti in qualsiasi momento della giornata era possibile contattarlo attraverso la piattaforma Discord.
\end{itemize}
\subsection*{Meeting settimanali}
É stato inoltre accordato con il tutor aziendale almeno un meeting a settimana in sede. Per questo motivo ogni lunedì è stato effettuato un aggiornamento con il tutor di quanto realizzato nella settimana; inoltre ad ogni meeting è stato riorganizzato il lavoro mancante in base alle tempistiche.
             % Processi
% !TEX encoding = UTF-8
% !TEX TS-program = pdflatex
% !TEX root = ../tesi.tex

%**************************************************************
\chapter{Protocolli di modellazione e trasferimento dati}
\label{cap:protocolli-trasmissione-dati}
%**************************************************************

\intro{Nel seguento capitolo vengono trattati dal punto di vista teorico i protocolli di modellazione e trasferimento dati, in particolare viene approfondito lo stile architetturale REST e il linguaggio di query GraphQL.
}\\

%**************************************************************
\section{Introduzione ai protocolli}
\subsection*{Modello architetturale client-server}
\label{client-server}
Prima di procedere nella spiegazione sui protocolli di trasferimento e modellazione dati, è necessario fare una breve introduzione sull'architettura delle Web Application moderne. Queste seguono ormai tutte un modello server - client, ovvero un modello architetturale che divide in due processi l'applicazione: un client che richiede servizi al server, il quale li esegue ritornando una risposta contenente l'esito dell'operazione. I protocolli di trasferimento e modellazione dati trovano il loro maggior utilizzo proprio nella comunicazione client - server, tuttavia prima di procedere con la loro spiegazione è necessario introdurre il concetto di Application Programming Interfaces.
\subsection*{Application Programming Interfaces}
 Comunemente dette API, ovvero l'acronimo di Application Programming Interfaces, sono interfacce comunemente realizzate per aggevolare la comunicazione tra  server e client. Ciascun applicativo/dispositivo è sviluppato con strutture di dati differenti che evolvono nel tempo, dunque risulta complessa la comunicazione tra queste entità. Le API giocano un ruolo fondamentale nello scambio di dati: infatti definiscono una interfaccia per la comunicazione, la quale è indipendente dall'implementazione specifica del dispositivo o dell'applicativo e permette di comunicare secondo delle regole specifiche riportate nella propria documentazione. Risultano dunque fondamentali nella comunicazione, collaborazione e integrazione di nuovi componenti applicativi.\\ \\
 Al giorno d'oggi le API vengono utilizzate dalla maggior parte delle web applications, dispositivi IoT, applicativi di vario genere e molto altro ancora.
 Nello specifico nella tesi si fa riferimento alle Web API, ovvero a quelle interfacce che sfruttano il protocollo HTTP per la comunicazione con altri applicativi/dispositivi. Ci sono diversi tipi di Web API, tra queste:
 \begin{itemize}
   \item \textbf{API pubbliche}: si tratta di API accessibili da tutti (possono essere anche a pagamento);
   \item \textbf{API private}: si tratta di API create con lo scopo di essere utilizzate solo ed esclusivamente all'interno dell'azienda;
   \item \textbf{API partner}: si tratta di API utilizzate tra aziende in collaborazione;
   \item \textbf{API composte}: si tratta di API differenti combinate tra loro per creare una sequenza di operazioni.
 \end{itemize}
La necessità di standardizzare il modo in cui vengono sviluppate le interfacce API ha portato dunque alla nascita dei protocolli sul trasporto di dati.
\subsection*{Protocolli di trasferimento dati}
Per protocolli di modellazione e trasferimento dati s'intende un insieme di regole, strutture e vincoli che regolano il funzionamento delle API. Permettono dunque di definire una sorta di standard al quale gli sviluppatori possono far riferimento per implementare e interagire con le API. Il termine protocollo non si addice perfettamente a tutte le varie tecnologie di data fetching, tuttavia a grandi linee può racchiuderle e dunque verrà utilizzato per questione di comodità.
\subsection*{I primi protocolli e loro evoluzione}
Al giorno d'oggi il protocollo di modellazione e trasferimento dati più utilizzato è sicuramente REST, tuttavia sono presenti anche altre tipologie di protocolli in utilizzo o che comunque sono state utilizzate in passato. Si tratta di tecnologie con lo stesso scopo, ma di natura completamente differente, di seguito le principali.
\subsubsection*{Remote Procedure Call}
Viene spesso indicato con l'acronimo RPC, si tratta di protocollo secondo il quale una procedura o subroutine viene invocata da un client esterno al server che deve eseguire la procedura, senza che il client conosca i dettagli del network. Viene utilizzato per chiamare processi in sistemi remoti, ma come fossero locali. \\
Di seguito riportata la definizione attribuita all'RPC dagli informatici Andrew Birrell e Bruce Nelson nel 1984:
  \begin{quoting}
    “Meccanismo sincrono che trasferisce il flusso di controllo e i dati attraverso una chiamata di procedura tra due spazi di indirizzo su una rete a banda stretta.”
  \end{quoting}
Come nelle chiamate a procedure locali, un RPC è una operazione sincrona che tiene in pausa il client fino al momento in cui ritorna il risultato della procedura invocata.
\subsubsection*{Simple Object Access Protocol}
Indicato spesso con l'acronimo SOAP, si tratta di un vero e proprio protocollo che definisce la struttura dei dati che devono essere trasferiti e come questi devono essere elaborati. Richiede esclusivamente il formato XML per trasferire dati e tipicamente viene utilizzato il protocollo HTTP per il trasferimento di file, tuttavia possono essere utilizzati anche protocolli differenti come ad esempio il protocollo SMTP. La struttura di un messaggio SOAP è composta da 3 principali componenti:
\begin{itemize}
  \item \textbf{Envelope}: necessario al fine di identificare il documento come messaggio SOAP;
  \item \textbf{Header}: è opzionale. Lo scopo dell'header nei messaggi SOAP è quello di trasportare indicazioni estranee al messaggio che si vuole trasportare, ma che vengono interpretate da i diversi nodi durante il cammino del messaggio;
  \item \textbf{Body}: il body contiene il vero e proprio messaggio che si vuole trasferire.
\end{itemize}
\subsubsection*{Ad hoc}
\section{Approfondimento sullo stile architetturale REST}
REST è l'acronimo di "Representational State Transfer" e si tratta di un tipo di stile architetturale introdotto da Roy Fielding nel 2000 e viene considerato al giorno d'oggi come uno standard per la realizzazione di web API. Si tratta di una astrazione degli elementi di un architettura di un sistema, del quale REST ne ignora i dettagli del'implementazione delle componenti e della sintassi del protocollo imponendo dei vincoli sul loro ruolo e sulla loro interazione.
\subsection*{Principi di un architettura REST}
\label{principi-REST}
Un architettura REST dunque deve rispettare alcuni principi, di seguito verranno elencati i sei gruppi definiti da Fielding.
\subsubsection*{Client-server}
Il primo principio sposa uno dei paradigmi cardine dell'informatica, ovvero il principio di \textit{Separation of concerns}, secondo il quale conviene sempre separare un sistema complesso in moduli distinti in modo che ognuno possa avere un proprio compito.\\
Ciò viene ripreso nell'architettura REST separando il client dal server, dunque dividendo due logiche diverse in due moduli distinti. Così facendo server e client possono essere implementati in maniera indipendente, usando qualsiasi lingua o tecnologia, basta che siano conformi al prossimo principio detto \textit{uniform interface}.
\subsubsection*{Uniform interface}
Si tratta di un principio fondamentale che differenzia le REST API da qualsiasi API non REST. Secondo questo prinpicio l'interazione tra componenti Web, dunque client, server e tutti gli intermediari del network, dipendono dalla uniformità delle loro interfacce.\\
I componenti Web dunque sono in grado di comunicare coerentemente seguendo quattro vincoli sull'interfaccia delineati da Fielding; questi sono:
\begin{itemize}
  \item \textbf{Identification of resources}: le risorse che vengono richieste devono essere identificate nella richiesta stessa, dunque specificandole nell'url;
  \item \textbf{Manipulation of resources through representations}: il client deve avere la rappresentazione delle risorse e deve poter sapere come modellarle sul server. L'idea alla base è che la rappresentazione (attraverso un qualsiasi formato, ad es. JSON, XML, ecc...) è una modo per interagire con le risorse, ma non è la risorsa stessa;
  \item \textbf{Self-descriptive messages}: in ciascun messaggio devono esser presenti le informazioni necessarie a descrivere come deve essere processata la richiesta;
  \item \textbf{Hypermedia as the Engine of Application State}: la rappresentazione dello stato di una risorsa deve includere i riferimenti alle risorse correlate. É dunque necessario includere i link per ciascuna risposta, così che il client possa navigare tra le altre risorse facilmente.
\end{itemize}
\subsubsection*{Layered System}
Secondo questo principio l'architettura di un applicativo deve essere composta da più strati. Ciascuno strato inoltre è cieco rispetto agli altri strati, tranne per quanto riguarda gli strati adiacenti. Questi layer possono essere composti da intermediari basati sul network i quali intercettano la comunicazione client-server con uno scopo specifico (ad esempio per questioni di sicurezza, caching, controllo del flusso dati, ecc...), possono essere ad esempio proxy e gateways. Per il principio di Layered System questi intermediari devono aderire alle interfacce al fine di mantenerne l'uniformità.
\subsection*{Cache}
Si tratta di uno dei vincoli fondamentali in un archiettura Web, secondo il quale un web server deve dichiarare la \textit{cacheability} di ciascuna risposta ritornata. Più specificatamente qualsiasi risposta di un server deve etichettare come cacheabili o meno i dati presenti all'interno di esso. Così facendo gli intermediari tra server e client e il client stesso sanno come comportarsi riguardo alla memorizzazione dei dati.
\subsection*{Stateless}
Il vincolo di stateless fa riferimento al fatto che un server non deve memorizzare lo stato dell'applicazione client. Questo implica però che ogni richiesta che il server riceve dal client deve essere sufficientemente dettagliata sullo stato del client affinché il server sia in grado di eseguirla. Dunque le richieste non sono correlate tra loro e per questo viene definito "stateless". \\
Questo vincolo porta un vantaggio fondamentale secondo il quale un server così facendo può gestire richieste da molti client. Può inoltre esser scalato molto più facilmente con l'aiuto ad esempio di un load balancer.
\subsection*{Code on demand}
Per ultimo troviamo il vincolo di code on demand, si tratta di un vincolo facoltativo secondo il quale la logica del client può essere aggiornata indipendentemente da quella lato server. Un esempio pratico lo troviamo nella signel web application le quali rispettano totalmente queste vincolo.

\section{Approfondimento sul linguaggio di query GraphQL}
\subsection*{Introduzione}
GraphQL è stato ideato da Facebook nel 2012 e condiviso e reso pubblico nel 2014. Al giorno d'oggi molte importanti applicazioni  utilizzano GraphQL, come ad esempio GitHub, Twitter, PayPal e Pinterest. Viene considerato come il principale competitor e possibile successore di REST nell'ambito del data fetching, tuttavia come verrà spiegato in seguito, oltre a svariati punti di forza e di innovazione ha anche alcuni problemi.\\
Più nello specifico GraphQL è un linguaggio di query per le APIs. Viene definito agnostico rispetto al mezzo di trasporto perché non dipende dal modo in cui client e server comunicano, ma solitamente viene utilizzato sul protocollo HTTP. Il principale punto di forza di GraphQL è la possibilità di specificare nella query esattamente i dati che si è interessati a ricevere, questo permette dunque di non occupare la rete per dati non richiesti. Altro importante punto di forza, ma che talvolta può risultare un problema, è che è fortemente tipizzato.\\
Affermare che GraphQL abbia lo scopo di servire esclusivamente come linguaggio di query può risultare riduttivo. Una dei principali motivi d'utilizzo di GraphQL è quello di riuscire a raggruppare tutti i dati e servizi di un'applicazione insieme in uno stesso posto, e fornire così un'interfaccia unica che risulti consistente, sicura e infine semplice da utilizzare.\\
GraphQL non specifica come deve essere costruita un'API, tuttavia ci sono cinque linee guida dette "Principi di desing" da tenere in considerazione durante lo sviluppo di un API:
\begin{itemize}
  \item \textbf{Hierarchical}: i tipi ricercati in una query GraphQL seguono una struttura gerarchica, infatti i tipi possono avere come campi altri tipi e così via. Inoltre i dati che vengono ritornati dalla query, vengono ritornati esattamente con la medesima struttura con cui sono stati richiesti;
  \item \textbf{Product centric}: le API sono inevitabilmente guidate dalle richieste dal client, per questo bisogna realizzarle in maniera flessibile cercando di tener conto delle richieste client per permettere quanto richiesto;
  \item \textbf{Strong typing}: un server GraphQL è supportato da un type system specifico a seconda dell'applicazione. Data una query, il server assicura che questa sia sintatticamente corretta, valida e che i tipi in gioco rispettino esattamente la struttura dei tipi definiti nel GraphQL schema;
  \item \textbf{Client-specified queries}: in GraphQL, la codifica della query avviene nel client e non nel server e si tratta di query che vanno a specificare campo per campo. Nella maggiorparte dei sistemi che non utilizzano GraphQL, il server determina quali dati ritornare. In GraphQL ciò non accade, vengono infatti ritornati solo i dati specificati dal client;
  \item \textbf{Introspective}: GraphQL è introspettivo, infatti i clients possono consultare a fondo il GraphQL schema e possono dunque vedere tutte le query disponibili, i vari tipi e i loro campi.
\end{itemize}
\subsection*{GraphQL schema}
GraphQL ha cambiato il modo di pensare alle APIs: queste non vengono più considerate come un insieme di endpoints dai quali ottenere dati ed eseguire servizi, ma vengono piuttosto considerate come una collezione di tipi e relazioni distribuiti su grafo.\\
Questa collezione di tipi e loro relazioni deve essere riportata in maniera chiara in una pagina chiamata \textbf{GraphQL Schema}. In questo schema devono essere definiti tutti quei tipi che si vuole esporre alle richieste del client per esser modificati, aggiunti, eliminati o semplicemente visualizzati. Il GraphQL schema viene realizzato con un linguaggio denominato SDL, acronimo di Schema Definition Language, il quale permette di specificare diverse tipologie di operazioni disponibili e i tipi, specificandone i campi e le relazioni tra di loro.\\
Durante la progettazione e la realizzazione di un server di backend possono essere seguite due differenti tecniche per la realizzazione dello schema e del resto del codice:
\begin{itemize}
  \item \textbf{Schema first}: questa tecnica si discosta dal classico ordine di realizzazione di un server di backend, infatti, come deducibile dal nome, prevede la realizzazione del GraphQL schema prima di realizzare i resolver delle API;
  \item \textbf{Conde first}: questa tecnica è opposta alla tecnica di Schema first e prevede la realizzazione prima dei resolver e solo in seguito dello schema.
\end{itemize}
Non è chiaro quale delle due tecniche sia migliore, sono aperti ancora diversi dibattiti a riguardo, in questo progetto è stato seguito il modello Schema first durante la migrazione, dunque è stato prima definito il GraphQL Schema e solo in seguito i resolver.\\
Grazie alla presenza del GraphQL Schema gli sviluppatori frontend sono in grado di conoscere nel dettaglio la struttura di ciascun tipo e delle varie query. Le APIs sviluppate con GraphQL si autodocumentano proprio perché è sufficiente la consultazione dello schema per comprendere la natura delle entità che si desidera interrogare o modellare.\\
\subsection*{Definizione dei tipi}
Come detto in precedenza la caratteristica principale di GraphQL è che si tratta di un linguaggio di query fortemente tipizzato. I tipi sono l'unità principale di un GraphQL schema.\\
Per tipo s'intende un oggetto costruito dettagliatamente campo per campo che deve poi corrispondere ad una entità nel backend dell'applicativo. Dunque all'interno dello schema dovranno essere definiti tutti i tipi che andranno a rappresentare la struttura dati dell'applicativo.\\
Un tipo può contenere come campi dati altri tipi che sono definiti nel medesimo schema. Segue un esempio di tipo dichiarato in uno schema GraphQL, in questo caso si tratta della dichiarazione di un Employee:
\begin{verbatim}
  type Employee {
    id: ID
    name: String!
    owns: Badge
    worksIn: Department
    worksOn: [Project]
    ...
  }
\end{verbatim}
In questo caso il tipo Employee avrà come campi:
\begin{itemize}
  \item \textbf{id}: un codice identificativo di tipo \textit{ID};
  \item \textbf{name}: un nome di tipo \textit{String}, il punto esclamativo indica che si tratta di un campo che non può essere nullo;
  \item \textbf{owns}: un badge di tipo \textit{Badge} per l'accesso al dipartimento;
  \item \textbf{worksIn}: un dipartimento di tipo \textit{Department} nel quale lavora l'impiegato,
  \item \textbf{worksOn}: una lista di progetti di tipo \textit{Project} al quale l'impiegato sta lavorando.
\end{itemize}
I tipi \textit{Department}, \textit{Badge} e \textit{Project} dovranno essere necessariamente definiti all'interno dello stesso GraphQL schema di Employee;
I built-in type che GraphQL mette a disposizione vengono detti \textbf{scalar type} e sono: \textit{Int}, \textit{String}, \textit{Boolean}, \textit{ID}, \textit{Float}.
È possibile inoltre dichiarare anche degli scalar type personalizzati con la keywork "\textit{scalar}", oppure delle enumerazioni attraverso l'utilizzo della keyword "\textit{enum}". È possibile infine unire diversi tipi, molto utile nel caso in cui si volesse ritornare uno tipo tra un insieme di tipi, questo è possibile farlo con la keyword "\textit{union}" come segue:
\begin{verbatim}
  union worker = Employee | Manager | Chief
\end{verbatim}
In questo caso sono stati uniti in un unico tipo \textit{worker} i tipi \textit{Employee}, \textit{Manager} e infine \textit{Chief}. Verrà molto utilizzato successivamente il tipo unione nella gestione degli errori di ritorno al client dopo l'esecuzione delle query.
\subsection*{Connessioni tra tipi}
GraphQL è cosi denominato perché oltre ad essere un Query Language come suggeriscono le ultime due lettere del nome, permette di definire connessioni di vario genere tra i tipi definiti nello schema, queste connessioni vanno di fatto a creare uno grafo composto da tipi interconnessioni, da questo deriva il prefisso \textit{Graph}.
È fondamentale durante la definizione del GraphQL schema riportare le relazioni nella maniera corretta delle entità corrispondenti nel database dell'applicativo. Un'ultima premessa prima di visualizzare i vari tipi di connessioni riguarda la direzionalità delle connessioni: in GraphQL risulta essere una buona pratica dare bidirezionalità alle connessioni ove possibile, questo con lo scopo di lasciare più flessibilità possibile allo sviluppatore client il quale dalla una query specifica può raggiungere diversi tipi e spostarsi nel grafo come più desidera.\\
Di seguito vengono elencate le varie relazioni con relativi esempi.
\subsubsection*{Connessione one-to-one}
Nelle relazioni one-to-one ad un tipo viene associata una sola istanza di un altro tipo e viceversa. Riprendendo il caso del tipo Employee riportato sopra, possiamo trovare una relazione del tipo one-to-one tra i tipi Employee e Badge. Di seguito la rappresentazione del grafo:
\begin{figure}[!h]
\centering
\includegraphics[width=0.5\linewidth]{immagini/one_to_one.pdf}
\caption{Connessione one-to-one.}
\label{one-to-one}
\end{figure}
\\ \\
Come mostrato in figura \ref{one-to-one} il collegamento tra i due tipi è definito come "owns" se si legge nel verso che parte da Employee per raggiungere Badge e fa riferimento all'omonimo campo di Employee. Se altrimenti la connessione si percorre nel verso opposto viene definita "isOwned", come l'omonimo campo di Badge, riportato in seguito:
\begin{verbatim}
  type Badge {
    id: ID
    isOwned: Employee
    ...
  }
\end{verbatim}
\subsubsection*{Connessione one-to-many}
In questo caso bisogna focalizzarsi sul campo \textit{worksIn} di Employee. Questo campo definisce la connessione con un elemento di tipo \textit{Department}, dunque a ciascun impiegato corrisponde un dipartimento nel quale lavora. Tuttavia pensando alla connessione in senso opposto a ciascun dipartimento possono corrispondere più impiegati. Segue dunque la rappresentazione della connessione nel grafo:\\

\begin{figure}[!h]
\centering
\includegraphics[width=0.4\linewidth]{immagini/one_to_many.pdf}
\caption{Connessione one-to-many.}
\label{one-to-many}
\end{figure}
\mbox{}\\
Come mostrato in figura \ref{one-to-many} il collegamento tra il dipartimento e i vari impiegati viene chiamato "employees" se si considera il verso che parte da Department per raggiungere Employee e corrisponde all'omonimo campo di Department. Tuttavia il collegamento è definito "worksIn" se la connessione viene percorsa nel verso opposto. Segue la rappresentazione del tipo Department:
\begin{verbatim}
  type Department {
    id: ID
    name: String!
    address: String!
    employees: [Employee]
  }
\end{verbatim}
\subsubsection*{Connessione many-to-many}
Consideriamo ora il campo \textit{worksOn} di Employee che collega ciascun impiegato con una lista di progetti ai quali sta lavorando. In questo caso però considerando il collegamento nel verso opposto anche ciascun progetto può avere più impiegati che ci lavorano. In questo caso si tratta di una relazione many-to-many e segue la rappresentazione nel grafo:
\begin{figure}[!h]
\centering
\includegraphics[width=0.4\linewidth]{immagini/many_to_many.pdf}
\caption{Connessione many-to-many.}
\label{many-to-many}
\end{figure}
\\ \\
% da sistemare questi doppi backslash
Le relazioni many-to-many non sono altro che l'unione di due relazioni one-to-many.
La connessione in questo caso, come nei casi precedenti, a seconda del verso in cui viene percorsa può esser definita come "worksOn" o "employees" come mostrato in figura \ref{many-to-many}. Di seguito la rappresentazione del tipo Project:
\begin{verbatim}
  type Project {
    id: ID
    name: String
    employees: [Employee]
  }
\end{verbatim}
\subsection*{Operazioni sui dati}
Come detto in precedenza GraphQL è un linguaggi di query e come tale permette di interrogare i dati o eseguire operazioni su di essi. Ci sono tre tipi di operazioni che possono essere fatte sui dati e queste sono: \textit{Query}, \textit{Mutation} e infine \textit{Subscription}.
\subsubsection{Query}
L'operazione di query viene utilizzata per richiedere dati da una determinata API ed equivale alla GET nel protocollo REST. È necessario dichiarare nel GraphQL schema la query che il programmatore backend desidera rendere disponibile, così facendo vengono dichiarati anche i tipi che possono eventualmente essere passati come argomenti e quelli che verranno ritornati dalla query.\\
Un esempio di dichiarazione di query che ritorna una lista di tutti gli oggetti di tipo Employee presenti in un determinato Department può essere:
\begin{verbatim}
  type query {
    employeesInDepartment(departmentId: ID!): [Employee]
  }
\end{verbatim}
In questo caso invocando la query \textit{employeesInDepartment} e passando come argomento alla query l'id (non può essere un valore nullo) del dipartimento, riceveremo come risposta un JSON contenente un campo "\textit{data}" contenente a sua volta la lista di impiegati, questo se la query ha avuto successo. In caso di insuccesso della query per qualsiasi motivo, ad esempio per avere passato un id errato, allora verrà ritornato un JSON con un campo "\textit{error}", contenente la descrizione dell'errore. \\
Dal punto di vista del client, qualora si volesse invocare questa query bisognerebbe strutturare la richiesta come segue:
\begin{verbatim}
  query {
    employeesInDepartment(id: "2BR4S") {
      id
      name
      worksOn {
        id
        name
      }
    }
  }
\end{verbatim}
Se la query non dovesse fallire, verrà ritornata una lista di impiegati e per ciascun impiegato verranno ritornati i campi specificati nella query dunque: l'\textit{id}, il \textit{name}, una lista di oggetti di tipo Project nel campo \textit{worksOn} per i quali bisognerà a loro volta specificare i campi ai quali si è interessati, in questo caso all' \textit{id} e al \textit{name} di ciascun progetto.\\
È inoltre possibile utilizzare gli argomenti delle query per controllare la quantità di dati che possono esser ritornati con un processo chiamato \textit{data paging}, oppure usarli per decidere in che ordine vogliamo che vengano ritornati i dati.
\subsubsection*{Mutation}
L'operazione di mutation viene utilizzata per eseguire modifiche sui dati. Equivale all'unione delle operazioni POST, DELETE, PUT, PATCH nel protocollo REST. Come per le query è necessario dichiarare le mutation che si voglion rendere disponibili al client.\\
Un esempio di mutation può essere:
\begin{verbatim}
  type mutation {
    addNewEmployee(employee: Employee!): Employee
  }
\end{verbatim}
In questo caso la mutation \textit{addNewEmployee} andrà ad aggiungere un nuovo impiegato nella struttura dati dell'applicativo. Il client per invocare questa mutation dovrà strutturare la richiesta come segue:
\begin{verbatim}
  mutation {
    addNewEmployee(employee: {
      name: "Mario"
    }){
      id
    }
  }
\end{verbatim}
In questa mutation è stato passato come argomento un oggetto di tipo Employee, del quale è stato specificato esclusivamente il nome (va obbligatoriamente specificato essendo un campo dichiarato non nullo). Come da definizione la mutation ritorna un oggetto di tipo Employee, del quale però in questo caso si vuole ricevere solo l'id generato. Come nel caso della query, se l'aggiunta dell'impiegato avrà successo, nel JSON di ritorno ci sarà un campo \textit{data} contenente l'id generato in seguito all' aggiunta dell'impiegato, in caso contrario sarà ritornato un JSON con un campo \textit{error} che descrive l'origine dell'errore.
\subsubsection*{Subscription}
L'ultimo tipo si chiama Subscription e si tratta di una funzione particolare resa disponibile in GraphQL, infatti grazie a questa funzione i client possono sottoscriversi ad una subscription e così facendo sarà il server ad inviare al client i dati richiesti non appena questi sono disponibili, dunque non è più necessario che sia il client a richiedere periodicamente i dati aggiornati.\\
Un esempio di definizione di una Subscription può essere:
\begin{verbatim}
  type subscription {
    newEmployeeAdded: Employee!
  }
\end{verbatim}
Il client che desidera sottoscriversi alla subscription \textit{newEmployeeAdded} dovrà mandare una richiesta strutturata come segue:
\begin{verbatim}
  subscription {
    newEmployeeAdded {
      id
      name
    }
  }
\end{verbatim}
Così facendo il server, appena viene aggiunto un nuovo impiegato, invierà direttamente al client i dati che il client ha specificato nella sottoscrizione, ovvero in questo caso l'\textit{id} e il \textit{name} dell'impiegato.\\
Essendo il server a dover inviare i dati al client e non il client che richiede i dati dal server, non è utilizzabile il protocollo HTTP per la comunicazione server - client, bisogna quindi utilizzare il protocollo WebSocket per aprire un canale di comunicazione a doppia via sopra un socket TCP.\\ \\
             % Kick-Off
% !TEX encoding = UTF-8
% !TEX TS-program = pdflatex
% !TEX root = ../tesi.tex

%**************************************************************
\chapter{Introduzione ai casi d'uso per l'analisi comparativa}
\label{casi-uso}
%**************************************************************

\intro{Il capitolo descrive in dettaglio il prototipo realizzato e la sua migrazione da REST a GraphQL. A seguire si illustrerà una ulteriore migrazione realizzata sull'applicativo aziendale SushiLab.}\\
%\subsection{Confronto con stakeholder}
\section{Prototipo}
\subsection{Introduzione}
Il prototipo è stato realizzato avendo presente i seguenti obbiettivi:
\begin{itemize}
  \item familiarizzare con le tecnologie Spring e Angular per la realizzazione rispettivamente di backend e frontend;
  \item familiarizzare con la realizzazione delle API sia con lo stile architetturale REST che con il linguaggio di query GraphQL;
  \item avere un caso d'uso ulteriore a conferma delle analisi che verranno poi ricavate dalla migrazione dell'applicativo SushiLab;
\end{itemize}
Si è scelto di realizzare un gestionale server - client in grado di gestire un insieme di impiegati. Ogni impiegato può lavorare in un dipartimento e può dividere il suo tempo su uno o più progetti. Una sede aziendale può contenere uno o più dipartimenti.
\subsection{Architettura ed entità}
\subsubsection*{Architettura generale dell'applicativo}
L'applicativo segue l' architettura raffigurata nell'immagine \ref{prototype-architecture}. Il backend del prototipo è stato realizzato con framework Spring Boot mentre il frontend con framework Angular.
\FloatBarrier
\begin{figure}[!h]
\centering
\includegraphics[width=1\linewidth]{immagini/prototypeArchitecture.pdf}
\caption{Architettura del prototipo}
\label{prototype-architecture}
\end{figure}
\FloatBarrier
Le richieste effettuate dai servizi del client vengono mappate sul REST controller del server il quale, dopo la rielaborazione interna delle richiesta, ritorna la risposta al client; quest'ultimo può dunque aggiornare la vista visualizzata dall'utente. In seguito vengono affrontati più nel dettaglio i moduli interni sia per quanto riguarda il server che il client.
\subsubsection*{Entità e relazioni}
\label{entity-relation}
Trattandosi di un gestionale aziendale semplificato si prevedono solo quattro entità principali:
\begin{itemize}
  \item \textbf{Employee}: l'impiegato che può lavorare ad uno o più progetti e in un dipartimento;
  \item \textbf{Project}: un progetto aziendale a cui partecipano più impiegati;
  \item \textbf{Site}: la sede aziendale, può avere più dipartimenti;
  \item \textbf{Department}: un dipartimento che appartiene ad una sede.
\end{itemize}
Ciascuna entità è caratterizzata dai campi presenti in figura \ref{ER-prototype}. I campi in grassetto rappresentano la chiave primaria, quelli sottolineati le chiavi esterne. Le relazioni presenti tra le varie entità sono:
\begin{itemize}
  \item \textbf{many to many}: è presente tra Employee e Project, infatti ciascun impiegato può lavorare a più progetti e ciascun progetto può avere più impiegati; viene scomposta in due relazioni \textit{one-to-many} con la tabella \textit{Employees\_Projects};
  \item \textbf{one to many}: è presente in due situazioni:
  \begin{itemize}
    \item tra Employee e Department, infatti ciascun impiegato deve lavorare almeno in un dipartimento, mentre ciascun dipartimento può ospitare più impiegati;
    \item tra Department e Site, ciascun dipartimento può appartenere esclusivamente ad una sede, mentre ciascuna sede può esser composta da più dipartimenti.
  \end{itemize}
\end{itemize}
%%% IMG
\FloatBarrier
\begin{figure}[!ht]
\centering
\includegraphics[width=1\linewidth]{immagini/ER_prototype.pdf}
\caption{Schema ER delle entità}
\label{ER-prototype}
\end{figure}
\FloatBarrier
%%% IMG
\subsection{Progettazione e realizzazione backend}
\subsubsection*{Architettura backend}
È stato deciso di seguire il pattern \textbf{controller - service - repository} per la realizzazione del server del prototipo.  Il pattern rispetta perfettamente il principio di "Separation Of Concerns". Di seguito in figura \ref{controller-service-repository} è possibile visualizzare l'architettura del backend.
%%%% IMG
\FloatBarrier
\begin{figure}[!h]
\centering
\includegraphics[width=0.8\linewidth]{immagini/controllerServiceRepository.pdf}
\caption{Architettura interna backend}
\label{controller-service-repository}
\end{figure}
\FloatBarrier
%%%% IMG
Nell'immagine \ref{sequence-diagram} è possibile visualizzare il diagramma di sequenza di una richiesta GET al server da parte del client. La richiesta in questo caso è andata a buon fine ritornando al client i dati richiesti.
\FloatBarrier
\begin{figure}[!h]
\centering
\includegraphics[width=1\linewidth]{immagini/getRequestFromClient.pdf}
\caption{Diagramma di sequenza di una richiesta GET elaborata dal server}
\label{sequence-diagram}
\end{figure}
\FloatBarrier
Il pattern prevedere la gestione delle entità e delle chiamate alle API attraverso tre strati:
\begin{itemize}
  \item \textbf{Controller layer}: è l'unico responsabile della comunicazione con i client, inoltre gestisce le interfacce REST e invoca lo strato di servizio;
  \item \textbf{Service layer}: è lo strato tra controller e repository, si occupa della business logic e qualora sia necessario visualizzare, salvare, modificare o eliminare dati allora comunica con lo strato di persistenza;
  \item \textbf{Repository layer}: si tratta dello strato inferiore dell'architettura, si occupa della gestione dei dati e delle loro modifiche nonché della gestione del database.
\end{itemize}
Durante la realizzazione viene seguito il percorso inverso rispetto a quanto visto nell'immagine \ref{controller-service-repository}: la realizzazione avviene partendo dallo strato di persistenza, dunque dalla creazione delle entità e del livello di repository.
\subsubsection*{Entità e repository}
A ciascuna entità del database viene fatta corrispondere una classe in Java quindi sono state realizzate quattro classi rappresentanti le entità Employee, Project, Department e Site.\\
Ciascuna classe implementa la classe \textit{Serializable}, così facendo è possibile serializzare i dati in flussi di byte. La serializzazione viene utilizzata poiché si tratta di dati che dovranno essere memorizzati nel database, è necessario serializzarli poiché abbandonano la Java Virtual Machine. Viene inoltre utilizzato un \textit{serialVersionUID} per attribuire una versione a ciascuna classe serializzabile in modo da garantire la consistenza delle classi memorizzate in banca dati.\\
Di seguito l'esempio dell'implementazione della classe Employee:
%%% IMG
\begin{lstlisting}[language=Java, title={Employee.java}, morecomment={[s][\color{DarkOrchid}]{\@}{\ }}, morecomment={[s][\color{OliveGreen}]{"}{"}},]
@Entity
@Table (name = "EMPLOYEE")
@Data
@NoArgsConstructor ()
@AllArgsConstructor ()
public class Employee implements Serializable {
    private static final long serialVersionUID = 8452515756703751450L;
    @Id
    @Column (name = "ID")
    private String id;
    private String name;
    private String surname;
    private Long salary;
    private Date birth;
    @ManyToMany
    @JoinTable (name = "EMPLOYEE_PROJECTS", joinColumns = { @JoinColumn (name = "EMPLOYEE_ID")},
            inverseJoinColumns = { @JoinColumn (name = "PROJECT_ID")})
    @JsonIgnore
    private Set<Project> projects = new HashSet<>();
    @ManyToOne
    @JsonBackReference(value = "employee_department")
    private Department department;
}
\end{lstlisting}
%\FloatBarrier
%\begin{figure}[!ht]
%\centering
%\includegraphics[width=0.9\linewidth]{immagini/employee_entity.png}
%\caption{Esempio di implementazione dell'entità Employee in Spring Boot.}
%\label{employee-entity-definition}
%\end{figure}
%\FloatBarrier
%%% IMG
Vengono attribuite alla classe Employee diverse annotazioni Spring, tra queste:
\begin{itemize}
  \item \textbf{@Entity}: si tratta dell'annotazione che permette di mappare la classe Employee come tabella nel database. Questa annotazione è resa disponibile dal modulo Spring Data JPA il quale implementa la specifica delle Java Persistence API attraverso Hibernate ORM e permette dunque il mapping della classi con corrispettive entità nel database;
  \item \textbf{@Table}: questa annotazione permette di specificare il nome della tabella generata o presente nel database durante la sua creazione o aggiornamento;
  \item \textbf{@Data}: grazie alla libreria \textit{lombok} è possibile, attribuendo alla classe questa annotazione, generare automaticamente tutti i metodi get e set per ciascun campo della classe;
  \item \textbf{@NoArgConstructor e @AllArgConstructor}: permettono di generare tutte le combinazioni di costruttori con e senza parametri.
\end{itemize}
Le annotazioni applicate agli attributi della classe sono:
\begin{itemize}
  \item \textbf{@Id}: associato al campo \textit{id}, specifica durante il mapping che si tratta della chiave primaria;
  \item \textbf{@GeneratedValue}: associato al campo \textit{id}, specifica che quando una nuova istanza dell'entità viene creata, deve essere generato un nuovo id randomico;
  \item \textbf{@ManyToMany} e \textbf{@ManyToOne}: queste specificano il tipo di relazione che è presente con le altre entità, sono rispettivamente associate ai campi \textit{projects}, con il quale Employee ha una relazione molti a molti e infine al campo \textit{department}, con il quale Employee ha una relazione molti a uno;
  \item \textbf{@JoinTable}: è associata al campo \textit{projects}, e poiché le relazioni molti a molti necessitano di una ulteriore tabella per la memorizzazione di tutte le associazioni, questa annotazione permette di specificarne il nome, ovvero \textit{EMPLOYEE\_PROJECTS} e i nomi delle due colonne, ovvero \textit{EMPLOYEE-ID} e \textit{PROJECT-ID};
  \item \textbf{@JsonIgnore}: associato al campo \textit{projects}, permette di escluderlo dalla serializzazione;
  \item \textbf{@JsonBackReference}: associato al campo \textit{department}, permette di dare una direzionalità alla relazione molti a uno con Department, fondamentale per evitare il problema della ricorsione infinita.
\end{itemize}
Analogamente sono state realizzate le classi corrispondenti alle entità Project, Department e Site.\\\\
A questo punto si procede con la realizzazione delle classi repository: ciascuna entità ha una propria repository corrispondente. Viene estesa l'interfaccia \textbf{JPArepository<T, ID>} con T il tipo della entità che si vuole gestire, mentre ID è il tipo della chiave primaria dell'entità T. La repository JPA estende le seguenti interfacce:
\begin{itemize}
  \item \textbf{CrudRepository<T, ID>}: contiene i metodi per gestire le classiche operazioni CRUD;
  \item \textbf{PagingAndSortingRepository<T, ID>}: contiene i metodi per gestire la paginazione e l'ordinamento;
\end{itemize}
Estendendo la JPARepository per ciascun tipo è possibile avere a disposizione diversi metodi per eseguire operazioni già implementate come: findAll, count, existById, SaveAndFlush, ecc...\\
Per realizzare ulteriori metodi è sufficiente rispettare le corretta sintassi del nome del metodo affinché Spring si occupi in automatico della sua implementazione.\\
Nel caso in cui si voglia realizzare una query personalizzata si utilizza l'annotazione \textbf{@Query} alla quale è possibile passare come attributo la query che desideriamo in linguaggio JPQL. Di seguito un esempio di quanto appena detto:
\begin{lstlisting}[language=Java, title={EmployeeRepository.java}, morecomment={[s][\color{DarkOrchid}]{@}{\ }},  morecomment={[s][\color{OliveGreen}]{"}{"}},]
@Repository ()
public interface EmployeeRepository extends JpaRepository<Employee, String> {
    List<Employee> findByName(String name);
    List<Employee> findBySurname(String surname);
    List<Employee> findByDepartment_id(Long id);
    @Query ("SELECT DISTINCT e FROM Employee e WHERE e.birth BETWEEN ?1 AND ?2")
    List<Employee> findByBornDateRange(@Param ("from") Date from, @Param ("to") Date to);
}
\end{lstlisting}
%%% IMG
%\FloatBarrier
%\begin{figure}[!ht]
%\centering
%\includegraphics[width=1\linewidth]{immagini/employee_repository.png}
%\caption{Esempio di implementazione della repository di Employee in Spring Boot.}
%\label{employee-repository}
%\end{figure}
%\FloatBarrier
%%% IMG
Oltre ai classici metodi di ricerca disponibili sono stati realizzati alcuni metodi per la ricerca di impiegati per nome, per cognome e per id di dipartimento in cui lavorano. Infine è stata realizzata una query personalizzata per la ricerca di impiegati nati in un range di date.\\
Da evidenziare l'uso dell'annotazione \textbf{@Repository} attribuita alla'interfaccia, fondamentale al fine di indicare che la classe fornisce meccanismi per modellare i dati del database.
\subsubsection*{Service}
Il compito dello strato di servizio è facilitare la comunicazione tra controller e repository oltre a realizzare la business logic dell'applicativo.\\
Per ciascun repository, dunque per ciascuna entità, è stato realizzato un servizio specifico per gestirne le logiche.\\
Al fine di rispettare i principi SOLID della programmazione, per questioni di loose coupling, IoC e semplicità nel testing, è stato scelto di implementare il pattern secondo il quale per ogni entità viene realizzata una interfaccia del corrispondente servizio e sua relativa implementazione. In figura \ref{service-serviceImpl} relativo all'entità Employee.
%%% IMG
\FloatBarrier
\begin{figure}[!ht]
\centering
\includegraphics[width=0.4\linewidth]{immagini/service_serviceImpl.pdf}
\caption{Esempio di implementazione dell'interfaccia EmployeeService}
\label{service-serviceImpl}
\end{figure}
\FloatBarrier
%%% IMG
A seguire viene riportato l'esempio dell'implementazione del servizio \textbf{EmployeeServiceImpl}:
\begin{lstlisting}[language=Java, title={EmployeeServiceImpl.java}, morecomment={[s][\color{DarkOrchid}]{@}{\ }},  morecomment={[s][\color{OliveGreen}]{"}{"}},]
@Service ()
public class EmployeeServiceImpl implements EmployeeService{
    @Autowired
    private EmployeeRepository employeeRepository;
    @Override
    public List<Employee> SelAll() { ... }
    @Override
    public Optional<Employee> SelById(String id){ ... }
    @Override
    public List<Employee> SelByName(String name){ ... }
    @Override
    public List<Employee> SelBySurname(String surname){ ... }
    @Override
    public List<Employee> SelByDepartment(Long id){ ... }
    @Override
    public List<Employee> SelByBornInRange(Date from, Date to){ ... }
    @Override
    public void InsEmployee(Employee employee){ ... }
    @Override
    public void DelEmployee(Employee employee){ ... }
}
\end{lstlisting}
%%% IMG
%\FloatBarrier
%\begin{figure}[!ht]
%\begin{mdframed}
%\centering
%\includegraphics[width=1\linewidth]{immagini/employeeServiceImpl.png}
%\end{mdframed}
%\caption{Classe EmployeeServiceImpl.}
%\label{employeeServiceImpl}
%\end{figure}
%\FloatBarrier
%%% IMG
In figura \ref{employeeServiceImpl} è possibile notare come la classe sia caratterizzata dall'annotazione \textbf{@Service} la quale viene utilizzata per indicare classi che contengono la business logic e per marcare la classe come service provider.\\\\
Le annotazioni utilizzate sono:
\begin{itemize}
  \item \textbf{@Autowired}: attribuita al campo \textit{employeeRepository} produce la dependency injection della repository \textit{EmployeeRepository} nel servizio;
  \item \textbf{@Override}: utilizzata da tutti i metodi, avverte che i metodi dichiarati nell'interfaccia vengono implementati;
\end{itemize}
Sono stati resi disponibili metodi semplici che vanno ad invocare, grazie alla dipendenza con la repository, le query già disponibili con la \textit{JPARepository<T, ID>} e la query vista precedenemente \textit{SelByBornInRage}.\\
Lo strato di servizio si incarica della gestione delle eccezioni estendendo il gestore  \textit{HandlerExceptionResolver}; in questo modo si migliora il comportamento di default che non sempre fornisce informazioni comprensibili sulla natura dell'errore.
\subsubsection*{Controller}
Lo strato di controller si occupa di gestire le richieste che il server riceve attraverso il protocollo HTTP e di mapparle inoltre sui metodi corrispondenti. Di seguito il codice del controller:
\begin{lstlisting}[title={EmployeeController.java}, language=Java, morecomment={[s][\color{DarkOrchid}]{@}{\ }},  morecomment={[s][\color{OliveGreen}]{"}{"}},]
@RestController
@RequestMapping (value = "api/employee")
public class EmployeeController {

    @Autowired ()
    private EmployeeService employeeService;

    @CrossOrigin
    @GetMapping (value = "/all", produces = "application/json")
    public ResponseEntity<List<Employee>> listAllEmployee(){ ... }

    @GetMapping (value = "name/{name}", produces = "application/json")
    public ResponseEntity<List<Employee>> listEmployeeByName(@PathVariable("name") String name){ ... }

    @GetMapping(value = "surname/{surname}", produces = "application/json")
    public ResponseEntity<List<Employee>> listEmployeeBySurname(@PathVariable("surname") String surname){ ... }

    @GetMapping (value = "department/{department_id}", produces = "application/json")
    public ResponseEntity<List<Employee>> listEmployeeByDepartment(@PathVariable("department_id") Long id){ ... }

    @GetMapping (value = "born/{from}/{to}", produces = "application/json")
    public ResponseEntity<List<Employee>> listEmployeeBornInDateRange(@PathVariable("from")
            @DateTimeFormat (pattern="yyyy-MM-dd") Date from, @PathVariable ("to") @DateTimeFormat  (pattern="yyyy-MM-dd") Date to){ ... }

    @PostMapping (value = "/insert", consumes = "application/json")
    public ResponseEntity<Employee> createEmployee(@RequestBody Employee employee){ ... }

    @DeleteMapping(value = "/delete/{id}")
    public ResponseEntity<?> deleteEmployee(@PathVariable ("id") String id){ ... }

    @PutMapping (value = "/update/{id}", consumes = "application/json")
    public ResponseEntity<?> updateEmployee(@PathVariable("id") String id, @RequestBody Employee employee){ ... }
}
\end{lstlisting}
%\FloatBarrier
%\begin{figure}[!ht]
%\begin{mdframed}
%\centering
%\includegraphics[width=1\linewidth]{immagini/EmployeeController.png}
%\end{mdframed}
%\caption{Classe EmployeeController.}
%\label{employeeController}
%\end{figure}
%\FloatBarrier
Questa classe permette di gestire e mappare le richieste a seconda dell'url dal quale proviene la richiesta. Di seguito le annotazioni associate alla classe:
\begin{itemize}
  \item \textbf{@RestController}: definisce la classe come controller di tipo REST;
  \item \textbf{@RequestMapping}: indica l'url al quale il client dovrà mandare le richieste per quello specifico controller.
\end{itemize}
Il controller ha una dipendenza con il servizio \textit{EmployeeService}, necessaria per l'invocazione dei suoi metodi. Ciascun metodo ritorna l'oggetto \textit{ResponseEntity<>} e consiste in una risposta HTTP alla quale viene modificato l'header e aggiunto il JSON nel body prima di essere ritornato. Le annotazioni associate ai metodi del controller sono:
\begin{itemize}
  \item \textbf{@GetMapping}: indica che si tratta di un metodo per la risoluzione di una richiesta GET; specifica l'url al quale ricevere la richiesta e ciò che viene ritornato, ovvero un file JSON;
  \item \textbf{@PostMapping}: indica che si tratta di un metodo per la risoluzione di una richiesta POST, specifica l'url al quale ricevere la richiesta e ciò che richiede in input, ovvero un file JSON passato attraverso il body della richiesta HTTP;
  \item  \textbf{@DeleteMapping}: indica che si tratta di un metodo per la risoluzione di una richiesta DELETE, specifica l'url al quale ricevere la richiesta;
  \item \textbf{@PutMapping}: indica che si tratta di un metodo per la risoluzione di una richieste PUT, specifica l'url al quale ricevere la richiesta e ciò che richiede in input, ovvero un file JSON passato attraverso il body della richiesta HTTP;
\end{itemize}
Oltre alle annotazioni sopra riportate sono presenti tra gli argomenti le annotazioni \textbf{@PathVariable} per indicare che si tratta di una variabile che verrà fornita nell'url nel posto definito dal nome specificato, \textbf{@RequestBody} ovvero un argomento che verrà fornito nel body della chiamata HTTP e infine \textbf{@DateTimeFormat} per specificare il formato del tipo di dato \textit{Date} che verrà passato dal client nell'url della richiesta.\\
Si rendono disponibili i seguenti metodi:
\begin{itemize}
  \item \textbf{listAllEmployee}: ritorna tutti gli impiegati presenti;
  \item \textbf{listEmployeeByName}: ritorna tutti gli impiegati con una determinato nome;
  \item \textbf{listEmployeeBySurname}: ritorna tutti gli impiegati con un determinato cognome;
  \item \textbf{listEmployeeByDepartment}: ritorna tutti gli impiegati di un dipartimento;
  \item \textbf{listEmployeeByBornInDataRange}: ritorna tutti gli impiegati nati in un determinato range di date;
  \item \textbf{updateEmployee}: consente l'aggiornamento dei campi dati di un impiegato;
  \item \textbf{createEmployee}: aggiunge un nuovo impiegato;
  \item \textbf{deleteEmployee}: rimuove un impiegato.
\end{itemize}
\subsubsection*{Testing API}
Essendo un prototipo incentrato sulla realizzazione delle API sono stati svolti in maniera semplice e veloce i test per gli strati di servizio e repository, essi non verranno riportati. Per quanto riguarda i test sui controller sono stati svolti test più approfonditi.\\
Per eseguire i test sui metodi del controller è stato scelto di utilizzare il framework JUnit5.\\
Il primo test realizzato è uno \textit{Smoke Test} che si incarica di verificare che il contesto Spring abbia effettivamente creato i controller dell'applicazione. In figura \ref{smoke-test} la classe \textit{SmokeTest} con una dipendenza per ciascun controller risolta con la dependency injection grazie all'annotazione \textbf{@Autowired}. Per ciascun controller è presente il relativo metodo di verifica, ciascuno annotato con \textbf{@Test} per indicare a JUnit5 che si tratta di un metodo di test.\\
Come possiamo notare è presente l'annotazione \textbf{@SpringBootTest}, necessaria per indicare a Spring Boot dove si trova la principale classe di configurazione e dunque avviare il contesto Spring. Da notare inoltre che ciascun metodo è dichiarato in maniera da poter sollevare eccezioni se necessario.
\begin{lstlisting}[language=Java, title={SmokeTest.java}, morecomment={[s][\color{DarkOrchid}]{@}{\ }},  morecomment={[s][\color{OliveGreen}]{"}{"}},]
@SpringBootTest (classes = PrototypeManagerApplication.class)
public class SmokeTest {
    @Autowired
    private EmployeeController employeeController;
    @Autowired
    private ProjectController projectController;
    @Autowired
    private DepartmentController departmentController;
    @Autowired
    private SiteController siteController;
    @Test
    public void contextLoadsEmployee() throws Exception {
        assertThat(employeeController).isNotNull();
    }
    @Test
    public void contextLoadsDepartment() throws Exception {
        assertThat(departmentController).isNotNull();
    }
    @Test
    public void contextLoadsProject() throws Exception {
        assertThat(projectController).isNotNull();
    }
    @Test
    public void contextLoadsSite() throws Exception {
        assertThat(siteController).isNotNull();
    }
}
\end{lstlisting}
%\FloatBarrier
%\begin{figure}[!ht]
%\begin{mdframed}
%\centering
%\includegraphics[width=0.7\linewidth]{immagini/SmokeTest.png}
%\end{mdframed}
%\caption{Classe \textit{SmokeTest} sulla creazione dei controller.}
%\label{smoke-test}
%\end{figure}
%\FloatBarrier
Passiamo ora ai test sui metodi del controller \textit{EmployeeController}. Di seguito è raffigurata la classe \textit{EmployeeControllerTest}, sono stati riportati solo i primi due metodi di test per questioni di spazio.
\begin{lstlisting}[language=Java, title={EmployeeControllerTest.java}, morecomment={[s][\color{DarkOrchid}]{@}{\ }},  morecomment={[s][\color{OliveGreen}]{"}{"}},]
@WebMvcTest (EmployeeController.class)
public class EmployeeControllerTest {
    @Autowired ()
    private MockMvc mockMvc;
    @Autowired ()
    private ObjectMapper objectMapper;
    @MockBean ()
    private EmployeeService employeeService;
    //Testing listAllEmployee()
    @Test ()
    void listAllEmployeeWihtNotEmptyList() throws Exception{
        int a = 0;
        int b = 0;
        int c = a+b;
    }
    @Test ()
    void listAllEmployeeWihtEmptyList() throws Exception{
        int a = 0;
        int b = 0;
        int c = a+b;
    }
}
\end{lstlisting}
La classe \textit{EmployeeControllerTest} utilizza l'annotazione \textbf{@WebMvcTest} per caricare nel contesto di test esclusivamente il controller testato.\\
Nella classe di test \textit{EmployeeControllerTest} sono dichiarate tre dipendenze fondamentali risolte con la dependency injection:
\begin{itemize}
  \item \textbf{MockMvc}: ogetto utilizzato per la simulazione di chiamate HTTP e la verifica della risposta;
  \item \textbf{ObjectMapper}: oggetto utilizzato nei test per la serializzazione e deserializzazione degli oggetti JSON in oggetti java e viceversa;
  \item \textbf{EmployeeService}: servizio che viene utilizzato nei test che usa l'annotazione \textbf{@MockBean} per aggiungere un mock dell'oggetto \textit{EmployeeService}.
\end{itemize}
Infine si riporta di seguito l'implementazione dettagliata del metodo di test \textit{createEmployee} seguendo il pattern \textit{Arrange - Act - Assert}.
\begin{lstlisting}[language=Java, title={createEmployee()}, morecomment={[s][\color{DarkOrchid}]{@}{\ }},  morecomment={[s][\color{OliveGreen}]{"}{"}},]
@Test ()
void createEmployee() throws Exception{
    //Arrange
    Employee employee = new Employee();
    employee.setName("Mario");
    employee.setSurname("Rossi");
    LocalDate date = LocalDate.of(1992, 05, 03);
    employee.setBirth(Date.from(date.atStartOfDay(ZoneId.systemDefault()).toInstant()));
    //Act & Assert
    mockMvc.perform(post("/api/employee/insert").contentType(MediaType.APPLICATION_JSON)
            .contentType(objectMapper.writeValueAsString(employee))).andExpect(status().isCreated())
            .andDo(print());
}
\end{lstlisting}
Nella prima parte (arrange) viene creato un nuovo Employee e gli viene assegnato un nome, un cognome e una data di nascita; successivamente (act) si esegue la chiamata GET utilizzando l'oggetto \textit{MockMvc} e passando nel body della richiesta HTTP l'oggetto Java Employee trasformato in JSON grazie all'\textit{ObjectMapper}. Per ultima la fase finale (assert) in cui si verifica che lo stato di ritorno sia \textit{created}, in caso contrario il test fallisce.
\subsection{Progettazione e realizzazione frontend}
Per quanto riguarda il frontend dell'applicativo è stato utilizzato il framework Angular per realizzare una Single-Page Application. Si tratta di un frontend minimale, realizzato con il solo scopo di comprendere e sviluppare la parte di comunicazione con il server utilizzando al massimo le REST API disponibili. Per questo motivo non è stata effettuata alcuna considerazione sull'utenza e sono stati tralascati completamente aspetti quali estetica grafica, accessibilità e responsive design.\\
\subsubsection*{Architettura}
Angular segue un'architettura ben precisa, in figura \ref{angular-architecture} è possibile visualizzarne la struttura:
\FloatBarrier
\begin{figure}[!ht]
\centering
\includegraphics[width=1\linewidth]{immagini/angularArchitecture.pdf}
\caption{Architettura della Web Application del prototipo.}
\label{angular-architecture}
\end{figure}
\FloatBarrier
Si tratta di una piattaforma modulare che basa la sua struttura sulla divisione in moduli. Ciascun modulo è composto da una componente ed un template detto anche vista; questi comunicano tra loro attraverso i metadati: l'utente che interagisce con la vista può attivare l'inizio di un procedura sulla componente. La componente spesso per l'esecuzione di alcune procedure invoca in suo aiuto i servizi iniettati nella componente stessa. \\\\
Nel corso della tesi quando si tratta il frontend verranno principalmente considerati i servizi; infatti quest'ultimi si incaricano della comunicazione con il server attraverso le API.
\subsubsection*{Funzionalità}
Di seguito si approfondisce la parte riguardante l'entità Employee, le altre entità sono state trattate in maniera analoga.\\
In linea con le query rese disponibili dal backend, devono essere presenti le seguenti funzionalità:
\begin{itemize}
  \item \textbf{Visualizzazione degli Employee}: devono poter essere visualizzata la lista degli impiegati permettendo la ricerca per: nome, cognome, data di nascita o dipartimento di appartenenza;
  \item \textbf{Aggiunta di un Employee}: deve essere possibile aggiungere un nuovo impiegato potendone specificare: nome, cognome, salario, data di nascita, dipartimento di appartenenza e infine i progetti che segue;
  \item \textbf{Aggiornamento di un Employee}: deve essere possibile aggiornare i campi dati di un impiegato specificandone l'id e il/gli campo/i da modificare;
  \item \textbf{Eliminazione di un Employee}: deve essere possibile eliminare un impiegato specificandone l'id.
\end{itemize}
Nel sezione sucessiva verrà mostrato come sono state implementate le funzionalità appena introdotte.
\subsubsection*{Realizzazione}
Inizialmente sono state create le interfacce per poter creare gli oggetti corrispondenti alle quattro entità del prototipo. A seguire l'esempio della dichiarazione della interfaccia dell'entità Employee:
\begin{lstlisting}[language=JavaScript, title={object.ts}]
export interface Employee {
  id: string;
  name: string;
  surname : string;
  salary: number;
  birth: Date;
  department: Department;
  project: [Project];
}
\end{lstlisting}
%\FloatBarrier
%\begin{figure}[!ht]
%\centering
%\includegraphics[width=0.3\linewidth]{immagini/employeeInterface.png}
%\caption{Interfaccia di Employee.}
%\label{employee-interface}
%\end{figure}
%\FloatBarrier
A seguire è stata implementata l'interaccia grafica inserendo nella pagina iniziale la possibilità di effettuare due scelte: la prima scelta riguarda l'entità, mentre la seconda riguarda l'operazione che si vuole eseguire sull'entità selezionata. In figura \ref{first-page} è possibile visualizzare le scelte:
\FloatBarrier
\begin{figure}[!ht]
\centering
\includegraphics[width=0.7\linewidth]{immagini/firstPage.png}
\caption{Prima pagina della Web Application.}
\label{first-page}
\end{figure}
\FloatBarrier
Una volta selezionate le due opzioni viene aggiunta dinamicamente la possibilità di specificare i dettagli della richiesta. Vengono di seguito mostrate le diverse interfacce che variano al variare della scelta selezionata e il servizio che permette di eseguire la chiamata alle API del backend.\\
In figura \ref{get-employee} abbiamo il caso in cui è stata selezionata l'entità Employee con operazione GET:
\FloatBarrier
\begin{figure}[!ht]
\centering
\includegraphics[width=0.7\linewidth]{immagini/getEmployee.png}
\caption{Opzioni disponibili con operazione GET selezionata.}
\label{get-employee}
\end{figure}
\FloatBarrier
Viene dunque esposta una lista contenente tutti gli impiegati, per ciascun impiegato vengono mostrati id, nome e cognome, salario e data di nascita. Successivamente viene visualizzata la scelta riguardante il campo per il quale eseguire la ricerca, in questo caso è stato scelto il nome, dunque inserendo nell'apposito campo di input il nome e selezionendo il bottone è possibile effettuare la ricerca per nome.\\
La scelta dell'opzione GET permette alla vista di invocare immediatamente il servizio che esegue una chiamata alle REST API del backend per ricevere la lista di tutti gli employee, stessa cosa accade anche quando premiamo il bottone "Esegui ricerca " dopo aver inserito il nome. Il servizio in questione si chiama \textit{EmployeeService} e di seguito ne è riportata l'implementazione:
\begin{lstlisting}[language=JavaScript, title={employee.service.ts}]
export class EmployeeService {

  url: String = 'http://localhost:8080/api/employee/';

  constructor(private http: HttpClient) { }

  allEmployees(){
    return this.http.get(this.url + "all")
    }
  employeeByName(name: String){ ... }
  employeeBySurname(surname: String){ ... }
  employeeByDepartment(id: String){ ... }
  employeeByBornInDateRange(from: Date, to: Date){ ... }
  createOrUpdateEmployee(employee: Employee){ ... }
  deleteEmployee(id: String){ ... }
}
\end{lstlisting}
%\FloatBarrier
%\begin{figure}[!ht]
%\centering
%\includegraphics[width=0.6\linewidth]{immagini/employeeService.png}
%\caption{Classe \textit{EmployeeService}.}
%\label{employee-service}
%\end{figure}
%\FloatBarrier
Il servizio realizzato per le chiamate alle REST API del backend permette di essere invocato direttamente dalla componenta relativa alla vista sulla quale è stata eseguita la selezione di un elemento HTML da parte dell'utente. È inoltre presente una dipendenza con l'oggetto \textit{HttpClient}, il quale viene iniettato nel servizio venendo passato come parametro al costruttore; la sua funzionalità è quella di permettere l'esecuzione delle chiamate alle API sul backend.\\\\
Proseguendo vengono visualizzati ora i casi in cui si sceglie come operazione quella di modificare o aggiungere un nuovo impiegato. In figura \ref{post-employee} è possibile visualizzare la porzione di interfaccia grafica per l'inserimento o l'aggiornamento di un impiegato:
\FloatBarrier
\begin{figure}[!ht]
\centering
\includegraphics[width=0.3\linewidth]{immagini/postEmployee.png}
\caption{Opzioni disponibili con operazione INSERT/UPDATE selezionata}
\label{post-employee}
\end{figure}
\FloatBarrier
Una volta inseriti tutti i valori per ciascun campo, tranne il campo ID che viene generato automaticamente, è possibile aggiungere un nuovo impiegato. Qualora invece si specificasse anche l'id, allora si tratterebbe di una modifica di un impiegato già esistente, questo solo se l'id inserito corrisponde veramente all'id di un employee. Anche in questo caso selezionando il bottone "Insert or update Employee!" viene invocato il metodo del servizio riportato in figura \ref{employee-service}.\\
Per finire in figura \ref{delete-employee} il caso in cui venga selezionata l'opzione DELETE:
\FloatBarrier
\begin{figure}[!ht]
\centering
\includegraphics[width=0.3\linewidth]{immagini/deleteEmployee.png}
\caption{Opzioni disponibili con operazione DELETE selezionata}
\label{delete-employee}
\end{figure}
\FloatBarrier
\subsubsection*{Testing}
I test di unità sono stati implementati seguendo il pattern \textit{Arrange - Act - Assert} sui metodi del servizio riportato in figura \ref{employee-service}. Di seguito è possibile visualizzare la funzione \textit{describe} la quale viene utilizzata per raggruppare un insieme di test:
\begin{lstlisting}[language=JavaScript, title={employee.service.spec.test}]
describe('EmployeeService', () => {
  let service: EmployeeService;
  let httpMock: HttpTestingController;

  beforeEach(() => {
    TestBed.configureTestingModule({
      imports: [HttpClientTestingModule],
      providers: [EmployeeService],
    });
    service = TestBed.inject(EmployeeService);
    httpMock = TestBed.inject(HttpTestingController);
  });
  afterEach(() => {
    httpMock.verify();
  });
  it('should retrieve all employees'), () => { ... }
  it('should retrive all employees by name'), () => { ... }
  it('should retrive all employees by surname'), () => { ... }
  it('should retrive all employees by department'), () => { ... }
  it('should create one employee'), () => { ... }
  it('should update one employee'), () => { ... }
  it('should delete one employee'), () => { ... }
\end{lstlisting}
Viene incluso il modulo Angular \textit{HttpClientTestModule}, necessario al fine di importare il servizio iniettabile \textit{HttpTestingController} utilizzato per il mocking e il flushing e dunque per eliminare le microtask in sospeso. Sono inoltre presenti le funzioni \textit{BeforeEach} utilizzata per configurare l'ambiente di test, dunque importare i moduli necessari e risolvere le dipendenze e la funzione \textit{afterEach} utilizzata per controllare, dopo l'esecuzion di ciascun test, che non siano rimaste richieste in sospeso. Infine, con la funzione \textit{it} viene definito un titolo per ciascun test e implementato il test stesso.\\
Più nello specifico si può visualizzare l'implementazione di un test per il metodo \textit{allEmployee} del servizio \textit{EmployeeService}:
\begin{lstlisting}[language=JavaScript, title={employee.service.spec.test}]
it('should retrieve all employees'), () => {
  //Arrange
  const employeesList: Employee[] = [
    {id: "F3BASD", name: "Mario", surname: "Rossi", salary: 1500, birth: new Date("05/04/1999")},
    {id: "AIF07S", name: "Giulio", surname: "Gialli", salary: 1500, birth: new Date("02/01/1988")},
    {id: "LOSF2D", name: "Serafino", surname: "Bianchi", salary: 1500, birth: new Date("22/11/1979")},
    {id: "PIFA2A", name: "Marta", surname: "Verdi", salary: 1500, birth: new Date("30/01/1995")},
  ];
  //Act & Assert
  service.allEmployees().subscribe(employees => {
    expect(employees).toBe(employeesList);
  })

  const request = httpMock.expectOne(`${service.url}/all`);
  request.flush(employeesList);
  expect(request.request.method).toBe('GET');

}
\end{lstlisting}
%\FloatBarrier
%\begin{figure}[!ht]
%\centering
%\includegraphics[width=0.8\linewidth]{immagini/allEmployeeTest.png}
%\caption{Test sul metodo \textit{allEmployee()} di \textit{EmployeeService}.}
%\label{get-all-test}
%\end{figure}
%\FloatBarrier
Viene inizialmente creato un array di Employee che ci si aspetterà di ricevere dalla chiamata e sucessivamente si verifica il risultato atteso tramite la funzione \textit{expect}.Per finire si controlla che l'url sia corretto, che venga effettuata una sola chiamata GET e che si tratti effettivamente di una richiesta GET.
\subsection{Migrazione da REST a GraphQL}
In questa sezione si affronta la migrazione sia del backend che del frontend da REST API a GraphQL API.
\subsubsection*{Migrazione Backend}
Il backend realizzato in Spring Boot con l'aiuto del modulo Spring Data REST per la realizzazione dei controller di REST API deve essere riscritto in parte utilizzando il modulo Spring GraphQL per la realizzazione dei controller in GraphQL.\\
É stata adottata la tecnica \textit{Schema First} che prevede la costruzione del GraphQL Schema ed, a seguire, l'implementazione dei resolver delle API. Nonostante i resolver fossero già realizzati per il prototipo in REST è comunque necessario adattarli all'utilizzo di questo nuovo protocollo.\\
Le quattro entità presenti, ovvero Employee, Project, Department e Site devono essere riportate tutte nel GraphQL schema poiché sono tutte accessibili dal client; esse saranno definite sia come entità di input che di output poiché GraphQL differenzia i due tipi nelle query, mutation e subscription. Di seguito è possibile visualizzare l'implementazione dei tipi \textit{Employee} e \textit{EmployeeInput}.
\begin{lstlisting}[language=GraphQL, title={schema.graphqls}]
type Employee {
  id: ID
  name: String
  surname: String
  salary: Int
  birth: Date
  projects: [Project]
  department: Department
}
\end{lstlisting}
\begin{lstlisting}[language=GraphQL, title={schema.graphqls}]
type EmployeeInput {
  id: ID
  name: String
  surname: String
  salary: Int
  birth: Date
  projectsId: [ID]
  departmentId: ID
}
\end{lstlisting}
Come si può notare i due tipi \textit{Employee} e \textit{EmployeeInput} non corrispondono completamente: i campi \textit{projects} e \textit{department} di \textit{Employee} non sono presenti in \textit{EmployeeInput}, o  meglio sono presenti ma in diversa forma. Il motivo di tale differenza si può sintetizzare in questo modo:
\begin{itemize}
  \item il tipo \textit{Employee} verrà ritornato al client su richiesta; dovrà contenere le varie istanze dei progetti a cui un impiegato sta partecipando come pure l'oggetto dipartimento;
  \item il tipo \textit{EmployeeInput} indica come fornire i dati necessari all'aggiunta/aggiornamento dell'impiegato; è sufficiente dunque passare gli id dei progetti coinvolti come pure l'id del dipartimento dell'impiegato.
\end{itemize}
A questo punto è possibile definire le query, le mutation ed eventualmente le subscription che si vogliono rendere disponibili al client. Di seguito è possibile visualizzare la dichiarazione delle query, mutation e della subscription riguardanti l'entità Employee.
\begin{lstlisting}[language=GraphQL, title={schema.graphqls}]
type Query {
  allEmployees: [Employee]
  employeeByName(name: String!): [Employee]
  employeeBySurname(surname: String!): [Employee]
  employeeByDepartment(id: ID!): [Employee]
  employeeBornInDateRange(from: Date!, to: Date!): [Employee]
}
type Mutation {
  addEmployee(employee: EmployeeInput!): Employee
  updateEmployee(employee: EmployeeInput!): Employee
  delEmployee(id: ID!): Employee
}
type Subscription {
  newEmployeeAdded():[Employee]
}
\end{lstlisting}
Il dato ritornato è di tipo \textit{Employee} per tutte le operazioni mentre quello di input è \textit{EmployeeInput}.\\
È stata aggiunta inoltre una subscription per sfruttare a pieno le funzionalità di GraphQL. Questa subscription permette al client di ricevere i nuovi impiegati aggiunti senza per forza richiederli.
\paragraph{Migrazione dei controller}
La migrazione non coinvolge gli strati di servizio e repository in quanto non è necessario apportare modifiche alle logiche e alla gestione del database. L'unico strato da riscrivere parzialmente è lo strato del controller.\\
Come prima cosa si realizzano le classi Java corrispondenti ai tipi di input; questo è necessario poiché nei metodi del controller è previsto come parametro di input la versione di input dell'impiegato.\\
L'implementazione del controller risulta differente rispetto a quella vista con le REST API. Di seguito è possibile visualizzarne la struttura ed i metodi:
\begin{lstlisting}[language=Java, title={EmployeeController.java}, morecomment={[s][\color{DarkOrchid}]{@}{\ }},  morecomment={[s][\color{OliveGreen}]{"}{"}},]
@Controller
public class EmployeeController {
    @Autowired
    private final EmployeeService employeeService;

    @QueryMapping
    Object employeeByName(@Argument String name) { ... }
    @QueryMapping
    Object employeeBySurname(@Argument String surname) { ... }
    @QueryMapping
    Object employeeByDepartment(@Argument String id) { ... }
    @QueryMapping
    Object employeeBornInDateRange(@Argument Date from, @Argument Date to) { ... }
    @MutationMapping
    Object addEmployee (@Argument EmployeeInput employee) { ... }
    @MutationMapping
    Object updateEmployee (@Argument EmployeeInput employee) { ... }
    @MutationMapping
    Object delEmployee(@Argument String id) { ... }
    @SubscriptionMapping
    Publisher<List<Employee>> newEmployeeAdded() { ... }
}
\end{lstlisting}
Come si può notare si tratta di un controller simile a quello REST nei metodi, tuttavia le annotazioni e alcune implementazioni differiscono. Vengono utilizzate le annotazioni:
\begin{itemize}
  \item \textbf{@Controller}: indica che si tratta di un generico controller (non più REST controller);
  \item \textbf{@QueryMapping}: attribuito ai metodi che risolvono le query;
  \item \textbf{@MutationMapping}: attribuito ai metodi che risolvono le mutation;
  \item \textbf{@SubscriptioMapping}: attribuito al metodo che risolve la subscription;
  \item \textbf{@Argument}: utilizzato per specificare gli argomenti.
\end{itemize}
Queste annotazioni sono fondamentali poiché permettono di effettuare il mapping tra i tipi, le query, le mutation e le subscription definiti nel GraphQL Schema e quelli definiti nel controller. Per questo motivo il nome di ciascun metodo o tipo negli argomenti deve corrispondere esattamente al nome nel GraphQL Schema, altrimenti il mapping non andrà a buon fine e verranno generati degli errori.\\
L'implementazione dei metodi del controller risulta simile a quella dei corrispettivi metodi nel REST controller. Tuttavia, trattandosi di una particolare funzionalità disponibile esclusivamente in GraphQL; viene riportata l'implementazione della subscription \begin{lstlisting}[language=Java, title={newEmployeeAdded()} morecomment={[s][\color{DarkOrchid}]{@}{\ }},  morecomment={[s][\color{OliveGreen}]{"}{"}},]
@SubscriptioMapping ()
Publisher<List<Employee>> newEmployeeAdded(){
  return subscriber -> Executors.newScheduledThreadPool(1).scheduleAtFixedRate(() -> {
    List<Employee> employees = employeeService.SelAll();
    subscriber.onNext(employees);
    }, 0, 2, TimeUnit.SECONDS);
}
\end{lstlisting}
L'implementazione della subscription permette al subscriber, in questo caso il client, di ricevere la lista aggiornata al verificarsi di un evento. L'oggetto \textit{Publisher<>} permette di inviare la lista aggiornata a tutti i subscribers in tempo reale, così facendo il client non dovrà richiedere periodicamente la lista con la query \textit{allEmployee}.
\paragraph{Strutture delle risposte HTTP in GraphQL e gestione degli errori}
Le risposte che il server GraphQL ritorna al client in seguito ad una query, mutation o subscription sono dei file JSON. Questi JSON possono contenere tre tipi di campi:
\begin{itemize}
  \item \textbf{Data}: ritornato nella risposta solo in caso di successo, assume valore \textit{null} al presentarsi di un errore;
  \item \textbf{Errors}: ritornato solo ed esclusivamente al verificarsi di uno o più errori, contiene le informazioni necessarie all'interpretazione dell'anomalia;
  \item \textbf{Extensions}: campo facoltativo a disposizione degli sviluppatori di backend, può essere utilizzato ad esempio per ritornare il timestamp dell'esecuzione della query.
\end{itemize}
Analizzando gli errori che si possono verificare durante l'esecuzione di una query, mutation o subscription nel server GraphQL, si individuano tre categorie:
\begin{itemize}
  \item \textbf{Errori di sintassi}: si tratta di errori nella sintassi della query. Se la query è sintatticamente scorretta l'errore sarà di tipo "Syntax Error"; in questo caso l'esecuzione della query non parte nemmeno;
  \item \textbf{Errori di validazione}: si verificano quando vengono specificati nella richiesta tipi non coerenti con lo schema. In questo caso l'errore si verifica durante l'esecuzione, dunque viene ritornato sia il campo \textit{data}, con valore null o con i dati parziali che il server è riuscito a recuperare, che il campo \textit{errors} che indica la presenza di uno o più errori di data fetching;
  \item \textbf{Errori nei resolver:} si verificano durante la risoluzione della query e non sono legati al fatto che la query sia formulata male o con tipi sbagliati. Questi errori devono essere gestiti lato server;
\end{itemize}
Quando si verificano errori di sintassi o validazione la risposta l'errore ritornato indica perfettamente la natura dell'errore. La terza tipologia di errore quasi sempre risulta incomprensibile dal client. Questo accade perché la gestione delle eccezioni durante l'esecuzione di una richiesta è presa in carico dall'interfaccia predefinita \textit{DataFetcherExceptionHandler} la quale permette la dichiarazione di più risolutori di eccezioni detti \textit{DataFetcherExceptionResolver}. Questi ultimi vengono invocati sequenzialmente fino a quando uno di essi è in grado di risolvere l'eccezione; in caso contrario viene ritornata una informazione di errore incomprensibile con categoria \textit{"INTERNAL\_ERROR"}. Questo per uno sviluppatore può essere un problema poiché non è facile individuare la provenienza dell'errore. Per questo motivo sono stati creati differenti resolvers per le varie eccezioni che possono verificarsi. Di seguito l'esempio di implementazione di un resolver che riscrive il metodo di risoluzione delle eccezioni \textit{resolveToSingleError} al fine di risolvere l'eccezione \textit{EmployeeNotFound}.
\begin{lstlisting}[language=Java, title={GraphQLExceptionResolver.java}, morecomment={[s][\color{DarkOrchid}]{@}{\ }},  morecomment={[s][\color{OliveGreen}]{"}{"}},]
@Component ()
public class GraphQLExceptionResolver extends DataFetcherExceptionResolverAdapter {
  @Override
  protected GraphQLError resolveToSingleError(Throwable ex, DataFetchingEnvironment env) {
    if(ex instanceof EmployeeNotFound) {
      return GraphQLErrorBuilder.newError()
                .errorType(ErrorType.NOT_FOUND)
                .message(ex.getMessage())
                .path(env.getExecutionStepInfo().getPath())
                .location(env.getField().getSourceLocation())
                .build();
    } else {
      return null;
    }
  }
}
\end{lstlisting}
Com'è possibile notare viene ritornato un oggetto di tipo \textit{GraphQLError} con i campi impostati correttamente. Al presentarsi dell'eccezione il client sarà ora in grado di comprendere il motivo di errore.
\subsubsection*{Migrazione Frontend}
La migrazione del frontend richiede uno sforzo minore rispetto al backend, infatti è sufficiente effettuare la migrazione esclusivamente dei servizi che si occupano dell'invocazione delle API.\\
Il servizio visualizzato precedentemente in figura \ref{employee-service} \textit{EmployeeService} non cambia la sua struttura; i metodi rimangono gli stessi, tuttavia l'implementazione del metodo deve cambiare. Oltre all'implementazione dei metodi cambia anche l'oggetto che permette il fetching dei dati dal server, infatti trattandosi ora di un GraphQL Server, sarà necessaria la dipendenza dall'oggetto \textit{Apollo} utilizzato nel fetching dei dati da un server GraphQL.\\
Si riporta l'implementazione del metodo  \textit{allEmployees}.
\begin{lstlisting}[language=JavaScript, title={allEmployees()}]
  allEmployees(){
      return this.apollo.watchQuery<any>({
        query: gql`
        {
          allEmployees {
            id
            name
            surname
            birth
            salary
          }
        }`,
        }).valueChanges;
      }
\end{lstlisting}
Attraverso il metodo \textit{watchQuery<>} dell'oggetto \textit{apollo} è possibile eseguire la query all'endopoint del server GraphQL. Il metodo ritorna un oggetto di tipo \textit{Observable} ed il subscriber dovrà occuparsi di gestire quanto ritornato dal servizio. Nel caso in cui la query sia andata a buon fine dovrà leggere le informazioni presenti nel campo \textit{"data"}; nel caso contrario dovrà occuparsi di mostrare il messaggio ed il tipo dell'errore presente nel campo \textit{"errors"}.
\section{SushiLab}
\label{sushi-lab}
\subsection{Confronto con stakeholder}
Circa dopo un mese dall'inizio dello stage è stato realizzato un confronto con gli stakeholder. Con il tutor interno dell'azienda SyncLab, ovvero Fabio Pallaro sono stati valutati alcuni applicativi già implementati che seguissero l'architettura server - client per effettuarne la migrazione da REST API a GraphQL API.\\
Tra tutti l'applicativo selezionato è stato \textbf{SushiLab}, il quale è stato realizzato da alcuni stagisti nei mesi precedenti. Il motivo che ci ha portato a concordare questa scelta è stato principalmente il fatto che fosse un applicativo con una struttura API articolata.
\subsection{Comprensione dell'applicativo e obbiettivi}
SushiLab è un applicativo che mira ad automatizzare il processo di tracciamento degli ordini nei ristoranti di sushi all you can eat. Si tratta di un applicativo composto da:
\begin{itemize}
  \item \textbf{Backend}: realizzato in Java con il framework Spring Boot con database PostgreSQL;
  \item \textbf{Frontend}: realizzato in Typescript con il framework Angular;
  \item \textbf{Componente di Machine Learning}: consente il riconoscimento dei piatti tramite la fotocamera del telefono;
  \item \textbf{Applicativo mobile}: la versione mobile della Web App.
\end{itemize}
Il mio intervento sulla Web Application prevede di lavorare sulla comunicazione tra backend e frontend, effettuando una migrazione completa delle REST API sviluppate trasformandole in GraphQL API, esattamente come è stato fatto per il prototipo.\\\\
Di seguito verranno analizzate le strutture del backend e del frontend tralasciando gli aspetti implementativi specifici già illustrati nel prototipo e concentrandosi principalmente sulla parte di comunicazione tramite API.
\subsection{Panoramica del backend}
Il backend sviluppato in Java con il framework Spring Boot segue l'architettura proposta per il prototipo, ovvero il pattern repository - service - controller. Di seguito verranno analizzate e descritte brevemente le componenti principali degli strati di persistenza e di business logic, si approfondirà maggiormente lo strato di controller.
\subsubsection{Entità}
In figura \ref{ER-sushilab} è possibile visualizzare lo schema ER delle entità presenti e le relazioni che intercorrono. Vengono tralasciati i campi per non entrare eccessivamente nel dettaglio.
\FloatBarrier
\begin{figure}[!ht]
\centering
\includegraphics[width=1\linewidth]{immagini/ER_sushilab.pdf}
\caption{Diagramma ER delle entità}
\label{ER-sushilab}
\end{figure}
\FloatBarrier
Sono dunque presenti le entità:
\begin{itemize}
  \item \textbf{PiattoUpload}: un tipo di piatto con allergeni e ingredienti;
  \item \textbf{PiattoPreview}: l'anteprima di un piatto associata ad un \textit{PiattoUpload};
  \item \textbf{SezionePreview}: un insieme di piatti di una particolare categoria;
  \item \textbf{Menù}: un insieme di sezioni di piatti, si tratta del menù di un ristorante;
  \item \textbf{Utente}: l'utente registrato, può anche essere un gestore e dunque essere il proprietario del ristorante e dunque poter apportare delle modifiche ai menù, ai piatti;
  \item \textbf{FasciaOraria}: una fascia oraria in cui può essere disponibile un Menù;
  \item \textbf{InformazioniPiatto}: le caratteristiche di un singolo utente riguardo ingredienti che non può mangiare, se gli è piaciuto un piatto specifico o se ha preferenze per particolari tipi di menù;
  \item \textbf{Blacklist}: l'insieme di ingredienti che un cliente preferisce non mangiare;
  \item \textbf{IdToken}: un token identificativo associato ad un utente dopo il login;
  \item \textbf{Session}: la situazione di un singolo tavolo in un determinato momento con determinati clienti, quando un cliente si siede ad un tavolo può aprire una sessione e invitare gli amici allo stesso tavolo;
  \item \textbf{Ordine}: gli ordini effettuati in una specifica sessione o da una specifica persona;
  \item \textbf{Allergeni}: le sostanze che tipicamente possono causare allergie;
  \item \textbf{Ingredienti}: i possibili ingredienti nei vari piatti.
\end{itemize}
Ciascuna di queste entità ha la corrispettiva classe Java nel backend. Per ciascuna entità è stata realizzata la corrispettiva repository estendendo l'interfaccia \textit{JPARepository<T, ID>} per la gestione dei dati e delle risorse nel database.
\subsubsection*{Servizi}
Lo strato di servizio presente prevede la definizione di solo due classi per la gestione delle logiche del backend. Di seguito vengono elencate le funzioni dei due servizi.
\paragraph{ClientService}
Il primo servizio gestisce le logiche delle possibili richieste del client. Questo servizio permette la gestione di:
\begin{itemize}
  \item \textbf{Utenti}: permette tutte le operazioni di login, registrazione di un nuovo utente, aggiunta dell'IdToken ad un utente all'apertura di una nuova sessione ad un tavolo e infine il recupero e modifica della password;
  \item \textbf{Preferenze Utente}: permette di modificare le preferenze di un utente e modificare, aggiungere o eliminare le valutazioni attribuite ai piatti; inoltre consente di visualizzare gli ingredienti nella blacklist di uno specifico utente;
  \item \textbf{Sessioni}: permette di creare ed eliminare le sessioni;
  \item \textbf{Ordini}: permette la gestione degli ordini da parte di un utente, di aggiungerne di nuovi o di visualizzarli;
  \item \textbf{Gestore}: permette di verificare che un utente corrisponda o meno al gestore;
  \item \textbf{Menu}: permette di visualizzare il menù di un ristorante.
\end{itemize}
\paragraph{GestoreService}
Il servizio \textit{GestoreService} gestisce le logiche riguardanti le operazioni realizzabili da un utente gestore, dunque permette:
\begin{itemize}
  \item \textbf{Piatti}: la gestione dei piatti, è possibile visualizzarli, eliminarli, modificarli e aggiungerne di nuovi;
  \item \textbf{Menu}: la gestione del menù, è possibile riorganizzarlo diversamente e attribuirgli differenti fasce orarie di disponibilità.
\end{itemize}
\subsubsection*{Controller}
Infine abbiamo lo strato più importante per gli scopi della tesi. Si tratta dello strato di controller il quale prevede cinque controller principali, di seguito verranno analizzati tutti i metodi di ciascun controller tralasciandone gli aspetti implementativi.
\paragraph{LoginController}
Il seguente controller dichiara una dipendenza con il servizio \textit{ClientService}. Di seguito la tabella dei metodi:
%//////////////////////// INIZIO TABELLA ////////////////////////
%//////////////////////// INIZIO TABELLA ////////////////////////
\begin{table}[!h]
\renewcommand{\arraystretch}{1.1}
\centering
\begin{tabular}{|c|c|c|l|}
\hline
\rowcolor[HTML]{C0C0C0}
{\color[HTML]{000000} \textbf{Metodo}} & {\color[HTML]{000000} \textbf{Parametri}} & {\color[HTML]{000000} \textbf{Descrizione}}                                           & {\color[HTML]{000000} \textbf{Azione}} \\ \hline
\textbf{eseguiLogin}                   & Utente utente             & \begin{tabular}[c]{@{}c@{}}Controlla ed effettua\\ il login di un utente\end{tabular} & POST                           \\ \hline
\end{tabular}
\caption{Metodi del controller \textit{LoginController}}
\end{table}
%//////////////////////// FINE TABELLA ////////////////////////
%//////////////////////// FINE TABELLA ////////////////////////
\paragraph{UtenteController}
Questo controller dichiara una dipendenza con il servizio \textit{ClientService} ed è composto dai metodi specificati nella seguente tabella:
%//////////////////////// INIZIO TABELLA ////////////////////////
%//////////////////////// INIZIO TABELLA ////////////////////////
\begin{table}[!h]
\renewcommand{\arraystretch}{1.2}
\begin{tabular}{|c|c|c|c|}
\hline
\rowcolor[HTML]{C0C0C0}
{\color[HTML]{000000} \textbf{Metodo}}        & {\color[HTML]{000000} \textbf{Parametri}}                                                                              & {\color[HTML]{000000} \textbf{Descrizione}}                                                                                                                & {\color[HTML]{000000} \textbf{Azione}} \\ \hline
{\color[HTML]{000000} \textbf{ottieniUtente}}          & {\color[HTML]{000000} String idPersona}                                                                                & {\color[HTML]{000000} \begin{tabular}[c]{@{}c@{}}Ricerca un utente\\ per id\end{tabular}}                                                                  & {\color[HTML]{000000} GET}             \\ \hline
\rowcolor[HTML]{EFEFEF}
{\color[HTML]{000000} \textbf{registraUtente}}         & {\color[HTML]{000000} Utente utente}                                                                                   & {\color[HTML]{000000} \begin{tabular}[c]{@{}c@{}}Aggiunge un nuovo\\ utente\end{tabular}}                                                                  & {\color[HTML]{000000} POST}            \\ \hline
{\color[HTML]{000000} \textbf{modificaStatoPreferiti}} & {\color[HTML]{000000} \begin{tabular}[c]{@{}c@{}}String idPersona\\ String idPiatto \\ String Param\end{tabular}}                   & {\color[HTML]{000000} \begin{tabular}[c]{@{}c@{}}Modifica i preferiti\\ di un utente,  con dei \\  nuovi piatti\end{tabular}}                                 & {\color[HTML]{000000} POST}            \\ \hline
\rowcolor[HTML]{EFEFEF}
{\color[HTML]{000000} \textbf{modificaValutazione}}    & {\color[HTML]{000000} \begin{tabular}[c]{@{}c@{}}String idPersona\\ String idPiatto\\ String param\end{tabular}} & {\color[HTML]{000000} \begin{tabular}[c]{@{}c@{}}Modifica la \\ valutazione di una \\ persona rispetto ad \\ un piatto con nuovi\\ parametri\end{tabular}} & {\color[HTML]{000000} POST}            \\ \hline
{\color[HTML]{000000} \textbf{recuperoPassword}}       & {\color[HTML]{000000} String param}                                                                                    & {\color[HTML]{000000} \begin{tabular}[c]{@{}c@{}}Invia un codice alla \\ mail del cliente per il\\ recupero password\end{tabular}}               & {\color[HTML]{000000} POST}            \\ \hline
\rowcolor[HTML]{EFEFEF}
{\color[HTML]{000000} \textbf{verificaCodice}}         & {\color[HTML]{000000} String param}                                                                                    & {\color[HTML]{000000} \begin{tabular}[c]{@{}c@{}}Verifica il codice\\ di recupero password\end{tabular}}                                                   & {\color[HTML]{000000} POST}            \\ \hline
{\color[HTML]{000000} \textbf{reimpostaPassword}}      & {\color[HTML]{000000} String param}                                                                                    & {\color[HTML]{000000} \begin{tabular}[c]{@{}c@{}}Imposta una nuova \\ password per un utente\end{tabular}}                                                 & {\color[HTML]{000000} POST}            \\ \hline
\rowcolor[HTML]{EFEFEF}
{\color[HTML]{000000} \textbf{aggiornaBlacklist}}      & {\color[HTML]{000000} \begin{tabular}[c]{@{}c@{}}String idPersona\\ Ingredienti ingredienti\end{tabular}}           & {\color[HTML]{000000} \begin{tabular}[c]{@{}c@{}}Aggiorna la blacklist \\ di una persona con i \\ nuovi ingredienti\end{tabular}}                          & {\color[HTML]{000000} POST}            \\ \hline
{\color[HTML]{000000} \textbf{ottieniBlacklist}}       & {\color[HTML]{000000} String idPersona}                                                                                & {\color[HTML]{000000} \begin{tabular}[c]{@{}c@{}}Ritorna la lista di \\ ingredienti nella \\ blacklist di un utente\end{tabular}}                          & {\color[HTML]{000000} GET}             \\ \hline
\end{tabular}
\caption{Metodi del controller \textit{UtenteController}}
\end{table}
%//////////////////////// FINE TABELLA ////////////////////////
%//////////////////////// FINE TABELLA ////////////////////////
\paragraph{TavoloController}
Anche il seguente controller dichiara una dipendenza con \textit{ClientService} ed i suoi metodi gestiscono gli utenti, gli ordini ed i tavoli.
% Please add the following required packages to your document preamble:
% \usepackage[table,xcdraw]{xcolor}
% If you use beamer only pass "xcolor=table" option, i.e. \documentclass[xcolor=table]{beamer}
% Please add the following required packages to your document preamble:
% \usepackage[table,xcdraw]{xcolor}
% If you use beamer only pass "xcolor=table" option, i.e. \documentclass[xcolor=table]{beamer}
\begin{table}[!h]
\renewcommand{\arraystretch}{1.2}
\begin{tabular}{|c|c|c|c|}
\hline
\rowcolor[HTML]{C0C0C0}
{\color[HTML]{000000} \textbf{Metodo}}                   & {\color[HTML]{000000} \textbf{Parametri}}                                                            & {\color[HTML]{000000} \textbf{Descrizione}}                                                                                                    & {\color[HTML]{000000} \textbf{Azione}} \\ \hline
{\color[HTML]{000000} \textbf{creaSessione}}             & {\color[HTML]{000000} String idPersona}                                                              & {\color[HTML]{000000} \begin{tabular}[c]{@{}c@{}}Permette la \\ creazione di una sessione\\  da parte di una persona\end{tabular}}             & {\color[HTML]{000000} POST}            \\ \hline
\rowcolor[HTML]{EFEFEF}
{\color[HTML]{000000} \textbf{ottieniSessione}}          & {\color[HTML]{000000} String idTavolo}                                                               & {\color[HTML]{000000} \begin{tabular}[c]{@{}c@{}}Ritorna la sessione \\ attiva in un determinato\\ tavolo\end{tabular}}                        & {\color[HTML]{000000} GET}             \\ \hline
{\color[HTML]{000000} \textbf{chiudiSessione}}           & {\color[HTML]{000000} String idTavolo}                                                               & {\color[HTML]{000000} \begin{tabular}[c]{@{}c@{}}Chiude la sessione in un\\ determinato tavolo\end{tabular}}                                   & {\color[HTML]{000000} DELETE}          \\ \hline
\rowcolor[HTML]{EFEFEF}
{\color[HTML]{000000} \textbf{ottieniOrdiniPersona}}     & {\color[HTML]{000000} \begin{tabular}[c]{@{}c@{}}String idTavolo\\ String idPersona\end{tabular}} & {\color[HTML]{000000} \begin{tabular}[c]{@{}c@{}}Ritorna gli ordini di \\ una persona in uno\\ specifico tavolo\end{tabular}}                  & {\color[HTML]{000000} GET}             \\ \hline
{\color[HTML]{000000} \textbf{modificaOrdiniPersona}}    & {\color[HTML]{000000} \begin{tabular}[c]{@{}c@{}}String idTavolo\\ String idPersona \\ ListaOrdini ordini\end{tabular}} & {\color[HTML]{000000} \begin{tabular}[c]{@{}c@{}}Modifica gli ordini di \\ una persona in uno \\ specifico tavolo\end{tabular}}                & {\color[HTML]{000000} POST}            \\ \hline
\rowcolor[HTML]{EFEFEF}
{\color[HTML]{000000} \textbf{ottieniOrdiniTavolo}}      & {\color[HTML]{000000} \begin{tabular}[c]{@{}c@{}}String idTavolo\\ String idPersona\end{tabular}} & {\color[HTML]{000000} \begin{tabular}[c]{@{}c@{}}Ritorna tutti gli ordini\\ effettuati dalle persone \\ di un determinato tavolo\end{tabular}} & {\color[HTML]{000000} GET}             \\ \hline
{\color[HTML]{000000} \textbf{ottieniGliOrdiniInArrivo}} & {\color[HTML]{000000} \begin{tabular}[c]{@{}c@{}}String idTavolo\\ String idPersona\end{tabular}} & {\color[HTML]{000000} \begin{tabular}[c]{@{}c@{}}Ritorna gli ordini che\\ sono quasi pronti per \\ essere consegnati\end{tabular}}             & {\color[HTML]{000000} GET}             \\ \hline
\rowcolor[HTML]{EFEFEF}
{\color[HTML]{000000} \textbf{spostaGliOrdiniInArrivo}}  & {\color[HTML]{000000} String idTavolo}                                                               & {\color[HTML]{000000} \begin{tabular}[c]{@{}c@{}}Imposta come "in arrivo"\\ gli ordini di un \\ determinato tavolo\end{tabular}}                 & {\color[HTML]{000000} POST}            \\ \hline
\end{tabular}
\caption{Metodi del controller \textit{TavoloController}}
\end{table}
%//////////////////////// FINE TABELLA ////////////////////////
%//////////////////////// FINE TABELLA ////////////////////////
%//////////////////////// INIZIO TABELLA ////////////////////////
%//////////////////////// INIZIO TABELLA ////////////////////////
\paragraph{MenuController}
Il \textit{MenuController} richiede la dipendenza con il servizio \textit{ClientService}. Ha i seguenti metodi:
\begin{table}[!h]
\centering
\renewcommand{\arraystretch}{1.2}
\begin{tabular}{|c|c|c|c|}
\hline
\rowcolor[HTML]{C0C0C0}
{\color[HTML]{000000} \textbf{Metodo}}                   & {\color[HTML]{000000} \textbf{Parametri}}                                                            & {\color[HTML]{000000} \textbf{Descrizione}}                                                                                                    & {\color[HTML]{000000} \textbf{Azione}} \\ \hline
{\color[HTML]{000000} \textbf{ottieniMenu}}              & {\color[HTML]{000000} \begin{tabular}[c]{@{}c@{}}int idMenu\\ String idPersona\end{tabular}}         & {\color[HTML]{000000} \begin{tabular}[c]{@{}c@{}}Ritorna il menu \\ ricercato per id\end{tabular}}                                             & {\color[HTML]{000000} GET}             \\ \hline
\rowcolor[HTML]{EFEFEF}
{\color[HTML]{000000} \textbf{ottieniListaFasce}}        & {\color[HTML]{000000} int idMenu}                                                                    & {\color[HTML]{000000} \begin{tabular}[c]{@{}c@{}}Ritorna le fasce orarie\\ di un certo menu\end{tabular}}                                      & {\color[HTML]{000000} GET}             \\ \hline
{\color[HTML]{000000} \textbf{ottieniListaPreferiti}}    & {\color[HTML]{000000} \begin{tabular}[c]{@{}c@{}}int idMenu\\ String idPersona\end{tabular}}         & {\color[HTML]{000000} \begin{tabular}[c]{@{}c@{}}Ritorna la lista dei\\ piatti preferiti di una \\ persona\end{tabular}}                       & {\color[HTML]{000000} GET}             \\ \hline
\end{tabular}
\caption{Metodi del controller \textit{MenuController}}
\end{table}
%//////////////////////// FINE TABELLA ////////////////////////
%//////////////////////// FINE TABELLA ////////////////////////
\paragraph{GestoreController}
Infine è presente il \textit{GestoreController} con dipendenza dal servizio \textit{GestoreService}. I metodi che lo caratterizzano sono:
%//////////////////////// INIZIO TABELLA ////////////////////////
%//////////////////////// INIZIO TABELLA ////////////////////////
% Please add the following required packages to your document preamble:
% \usepackage[table,xcdraw]{xcolor}
% If you use beamer only pass "xcolor=table" option, i.e. \documentclass[xcolor=table]{beamer}
\begin{table}[!h]
\renewcommand{\arraystretch}{1.2}
\begin{tabular}{|c|c|c|c|}
\hline
\rowcolor[HTML]{C0C0C0}
{\color[HTML]{000000} \textbf{Metodo}}             & {\color[HTML]{000000} \textbf{Parametri}}                                                                               & {\color[HTML]{000000} \textbf{Descrizione}}                                                                          & {\color[HTML]{000000} \textbf{Azione}} \\ \hline
{\color[HTML]{000000} \textbf{ottieniListaPiatti}} & {\color[HTML]{000000} String idPersona}                                                                                 & {\color[HTML]{000000} \begin{tabular}[c]{@{}c@{}}Ritorna la lista di \\ piatti di una persona\end{tabular}}          & {\color[HTML]{000000} GET}             \\ \hline
\rowcolor[HTML]{EFEFEF}
{\color[HTML]{000000} \textbf{nuovoPiatto}}        & {\color[HTML]{000000} \begin{tabular}[c]{@{}c@{}}PiattoUpload piatto\\ String idPersona\end{tabular}}                   & {\color[HTML]{000000} \begin{tabular}[c]{@{}c@{}}Aggiunge un nuovo\\ piatto\end{tabular}}                            & {\color[HTML]{000000} POST}            \\ \hline
{\color[HTML]{000000} \textbf{ottieniPiatto}}      & {\color[HTML]{000000} String idPiatto}                                                                                  & {\color[HTML]{000000} \begin{tabular}[c]{@{}c@{}}Ritorna un piatto\\ per id\end{tabular}}                            & {\color[HTML]{000000} POST}            \\ \hline
\rowcolor[HTML]{EFEFEF}
{\color[HTML]{000000} \textbf{eliminaPiatto}}      & {\color[HTML]{000000} \begin{tabular}[c]{@{}c@{}}String idPersona\\ String idPiatto\end{tabular}}                       & {\color[HTML]{000000} \begin{tabular}[c]{@{}c@{}}Elimina un piatto\\ per id\end{tabular}}                            & {\color[HTML]{000000} DELETE}          \\ \hline
{\color[HTML]{000000} \textbf{aggiornaPiatto}}     & {\color[HTML]{000000} \begin{tabular}[c]{@{}c@{}}String idPiatto\\ String idPersona\\ PiattoUpload piatto\end{tabular}} & {\color[HTML]{000000} \begin{tabular}[c]{@{}c@{}}Aggiorna un nuovo\\ piatto con i campi \\ specificati\end{tabular}} & {\color[HTML]{000000} PUT}             \\ \hline
\rowcolor[HTML]{EFEFEF}
{\color[HTML]{000000} \textbf{ottieniListaMenu}}   & {\color[HTML]{000000} String idPersona}                                                                                 & {\color[HTML]{000000} Ritorna la lista dei menu}                                                                     & {\color[HTML]{000000} GET}             \\ \hline
{\color[HTML]{000000} \textbf{nuovoMenu}}          & {\color[HTML]{000000} \begin{tabular}[c]{@{}c@{}}String idPersona\\ Menu menu\end{tabular}}                             & {\color[HTML]{000000} \begin{tabular}[c]{@{}c@{}}Aggiunge un nuovo\\ menu\end{tabular}}                              & {\color[HTML]{000000} POST}            \\ \hline
\rowcolor[HTML]{EFEFEF}
{\color[HTML]{000000} \textbf{ottieniMenu}}        & {\color[HTML]{000000} \begin{tabular}[c]{@{}c@{}}String idPersona\\ String idMenu\end{tabular}}                         & {\color[HTML]{000000} Ritorna un menu per id}                                                                        & {\color[HTML]{000000} GET}             \\ \hline
{\color[HTML]{000000} \textbf{aggiornaMenu}}       & {\color[HTML]{000000} \begin{tabular}[c]{@{}c@{}}String idPersona\\ String idMenu\\ Menu menu\end{tabular}}             & {\color[HTML]{000000} \begin{tabular}[c]{@{}c@{}}Aggiorna un \\ determinato menu\end{tabular}}                       & {\color[HTML]{000000} PUT}             \\ \hline
\rowcolor[HTML]{EFEFEF}
{\color[HTML]{000000} \textbf{eliminaMenu}}        & {\color[HTML]{000000} \begin{tabular}[c]{@{}c@{}}String idPersona\\ String idMenu\end{tabular}}                         & {\color[HTML]{000000} \begin{tabular}[c]{@{}c@{}}Elimina un \\ determinato menu\end{tabular}}                        & {\color[HTML]{000000} DELETE}          \\ \hline
\end{tabular}
\caption{Metodi del controller \textit{GestoreController}}
\end{table}
%//////////////////////// FINE TABELLA ////////////////////////
%//////////////////////// FINE TABELLA ////////////////////////
\newpage
\subsubsection{Panoramica del frontend}
\paragraph{Architettura}
L'applicativo è stato realizzato in Angular e dunque è composto da:
\begin{itemize}
  \item \textbf{Moduli}: utili per la separazione delle logiche; sono presenti diversi moduli come: il menu di navigazione, una sezione ordini, una per la guida, una per il gestore ed altri ancora;
  \item \textbf{Servizi}: sono presenti diversi servizi per la gestione della comunicazione con il backend, questo è lo strato in cui avverrà la migrazione;
  \item \textbf{Interceptor}: necessari per intercettare le chiamate HTTP e gestire il token di autenticazione;
  \item \textbf{Guardie}: permettono di bloccare l'accesso ad alcune funzionalità secondo diverse logiche, come ad esempio evitare di mostrare ad un utente che non è un gestore la pagina di gestione del menù;
  \item \textbf{Components}: previsti da Angular sono necessari per gestire le viste.
\end{itemize}
\subsubsection{Casi d'uso o Funzionalità}
Nella seguente sezione verranno illustrati brevemente i principali attori e i loro casi d'uso. Sono presenti due attori principali:
\begin{itemize}
  \item \textbf{Cliente}: si tratta del cliente del ristorante; a sua volta si suddivide in:
    \begin{itemize}
      \item \textit{Cliente non riconosciuto}: cliente non loggato; può effettuare tutte le azioni di consultazione del menu e dei piatti, ordinare i piatti, unirsi ad un tavolo, generare una sessione, ricercare i ristoranti e infine tutte le azioni riguardanti il login o registrazione;
      \item \textit{Cliente riconosciuto}: cliente loggato; estende il cliente non riconosciuto; potrà gestire i preferiti, comporre la propria blacklist, valutare i piatti o eseguire il logout;
    \end{itemize}
  \item \textbf{Gestore}: rappresenta il gestore di un ristorante; a sua volta si suddivide in:
    \begin{itemize}
      \item \textit{Gestore non riconosciuto}: gestore non loggato; può registrarsi o effettuare il login;
      \item \textit{Gestore riconosciuto}: gestore loggato con account da gestore; può effettuare il logout, visualizzare i clienti e i loro dati, attribuire punti o sottrarne ai clienti, confermare gli ordini, modificare e aggiungere piatti, menù e relative fasce orarie;
      \item \textit{Cameriere}: cameriere del ristorante; può visualizzare le ordinazioni di un determinato tavolo;
    \end{itemize}
\end{itemize}
\subsection{Migrazione del backend}
La migrazione del backend avviene secondo la tecnica \textit{Schema First}. È stata suddivisa in tre parti principali:
\begin{itemize}
  \item \textbf{Analisi sui tipi}: in questa prima fase si analizzano i tipi che devono essere esposti al client, dunque si procede con la scrittura del GraphQL Schema definendo i tipi;
  \item \textbf{Analisi sulle API}: sono state analizzate le API che si vuole migrare e successivamente sono state dichiarate nel GraphQL Schema;
  \item \textbf{Migrazione dei controller}: è stato adattato il codice del controller a GraphQL;
\end{itemize}
\subsubsection*{Analisi sui tipi e loro definizione}
É necessario analizzare e dunque riportare solo le entità che devono essere esposte al client. Oltre a queste è necessario definire nuovi tipi a seconda delle necessità. Di seguito vengono esposti i tipi che è stato scelto di dichiarare nel GraphQL Schema:
\begin{itemize}
  \item \textbf{Tipi di output}: sono i tipi che verranno ritornati al client, tra questi troviamo:
  \begin{itemize}
    \item \textit{Menu}: per ritornare il menu;
    \item \textit{FasciaOraria}: per ritornare le fasce orarie riguardo la disponibilità di un menù;
    \item \textit{Piatto}: per ritornare uno o più piatti; nella versione erano presenti due diverse rappresentazioni di piatto, qui ne sono stati uniti i campi sotto un unico tipo, così facendo backend e frontend potranno decidere quali campi sfruttare e quali no a seconda dell'utilizzo che deve essere fatto di questo oggetto;
    \item \textit{PiattoPreferito}: si differenzia dal tipo \textit{Piatto} in quanto deve specificare le preferenze dell'utente e le sue valutazioni riguardo ad uno specifico piatto;
    \item \textit{LoginSession}: per ritornare i dati di login, comprende una stringa per il token della sessione di login e la durata in secondi;
    \item \textit{Ordine}: si tratta di un oggetto che specifica un piatto, la quantità ordinata ed infine una eventuale nota;
    \item \textit{Utente}: per creare un nuovo utente o effettuare il login;
    \item \textit{Blacklist}: per ritornare la blacklist di un utente;
    \item \textit{Sessione}: per ritornare la sessione attiva in un determinato tavolo;
  \end{itemize}
  \item \textbf{Tipi di input}: sono i tipi che verranno forniti in input dal client, tra questi troviamo:
  \begin{itemize}
    \item \textit{PiattoInput}: per aggiunta e modifica di uno o più piatti che si ricevono in input dal client;
    \item \textit{MenuInput}: per aggiunta e modifica di uno o più menu;
    \item \textit{OrdineInput}: per aggiunta di un ordine, specifica il piatto, la sua quantità ed eventuali note da parte del cliente;
    \item \textit{UtenteInput}: per creare un nuovo utente o per effettuare il login;
  \end{itemize}
\end{itemize}
I tipi elencati sono esclusivamente quelli necessari per la definizione delle query, mutation e subscription. Verrà poi realizzato il tipo \textit{OperationSuccess} che viene ritornato, in caso di successo, per tutte quelle mutation che eseguono un'operazione e non ritornano dati. Questo tipo è composto da un solo campo String contenente un eventuale messaggio.\\ \\
A questo punto sono state create in Java le classi corrispondenti a tutti i tipi di input. Per i tipi di output questa operazione non è necessaria poiché i loro campi corrispondono alle entità in Java. Fanno eccezione i tipi \textit{Piatto}, \textit{PiattoPreferito} e \textit{LoginSession} poiché si differenziano leggermente dalle classi Java corrispondenti.
\subsubsection*{Analisi sulle API e loro definizione}
Questa fase prevede la dichiarazione delle query, mutation e subscription nel GraphQL Schema. Si tratta di una fase fondamentale in quanto permette di definire i tipi di input e di output di ciascuna operazione possibile. Rispetto alla progettazione iniziale si è deciso di aggiungere delle subscription. Di seguito è possibile visualizzare le query, le mutation e le subscription per ciascuno dei cinque controller presenti.\\
Le query rese disponibili sono:
\begin{lstlisting}[language=GraphQL, title={schema.graphqls}]
type Query {
  ##############UtenteController##############
  ottieniUtente(idPersona: String!): Utente
  ottieniBlacklist(idPersona: String!): blacklist
  ##############TavoloController##############
  ottieniSessione(idTavolo: String!): Sessione
  ottieniOrdiniPersona(idTavolo: String!, idPersona: String!): [Ordine]
  ottieniOrdiniTavolo(idTavolo: String!, idPersona: String!): [Ordine]
  ottieniGliOrdiniInArrivo(idTavolo: String!, idPersona: String!): [Ordine]
  ##############MenuController##############
  ottieniMenu(idMenu: String!, idPersona: String!): Menu
  ottieniListaFasce(idMenu: String!): [FasciaOraria]
  ottieniListaPreferiti(idMenu: String!, idPersona: String!): [Piatto]
  ##############GestoreController##############
  ottieniListaPiatti(idPersona: String!): [Piatto]
  ottieniListaMenu(idPersona: String!): [Menu]
  ottieniMenu(idPersona: String!, idMenu: String!): Menu
}
\end{lstlisting}
Le mutation realizzate sono:
\begin{lstlisting}[language=GraphQL, title={schema.graphqls}]
type Mutation {
  ##############LoginController##############
  eseguiLogin(utente: UtenteInput!): LoginSession
  ##############UtenteController##############
  registraUtente(utente: UtenteInput!): OperationSuccess
  modificaStatoPreferiti(idPersona: String!, idPiatto: String!, piatti: [PiattoInput!]!): OperationSuccess
  modificaValutazione(idPersona: String!, idPiatto: String!, piatti: PiattoPreferito!): OperationSuccess
  recuperoPassword(mail: String!): OperationSuccess
  verificaCodice(codice: String!): OperationSuccess
  reimpostaPassword(password: String!): OperationSuccess
  aggiornaBlackList(idPersona: String!, ingredienti: [String!]!): OperationSuccess
  ##############TavoloController##############
  creaSessione(idPersona: String!): Sessione
  chiudiSessione(idTavolo: String!): OperationSuccess
  modificaOrdiniPersona(idPersona: String!, idTavolo: String!, ordini: [OrdineInput!]!): OperationSuccess
  spostaGliOrdiniInArrivo(idTavolo: String!): OperationSuccess
  ##############GestoreController##############
  nuovoPiatto(idPersona: String!, piatto: PiattoInput!): Piatto
  eliminaPiatto(idPersona: String!, idPiatto: String!): OperationSuccess
  aggiornaPiatto(idPersona: String!, idPiatto: String!, piatto: PiattoInput!): OperationSuccess
  nuovoMenu(idPersona: String!, menu: MenuInput!): Menu
  aggiornaMenu(idPersona: String!, menu: MenuInput!): OperationSuccess
  eliminaMenu(idPersona: String!, idMenu: String!): OperationSuccess
}
\end{lstlisting}
Sono state introdotte infine le nuove subscription:
\begin{lstlisting}[language=GraphQL, title={schema.graphqls}]
type Subscription {
  ordiniTavoliRealTime(): [Ordine]
  visualizzaSezioniRealTime(): [Sessione]
}
\end{lstlisting}
Queste subscription permettono al client di ricevere in tempo reale i nuovi ordini effettuati in qualsiasi tavolo e quando vengono occupati nuovi tavoli.
\subsubsection*{Migrazione dei controller}
L'ultimo step per completare la migrazione del backend prevede la riscrittura dei controller. Per effettuare questo passaggio sono stati seguiti i seguenti punti:
\begin{itemize}
  \item \textbf{Cambio tipi}: sono stati cambiati i tipi negli argomenti, negli oggetti di ritorno e nell'implementazione dei vari resolver; questo perchè sono stati introdotti alcuni nuovi tipi ed è dunque necessario rispettare il mapping tra le dichiarazioni delle query, mutation e subscription nel GraphQL Schema e i corrispettivi resolver nei controller;
  \item \textbf{Cambio annotazioni}: sono state sostituite le annotazioni REST con quelle GraphQL;
  \item \textbf{Cambio struttura controller}: sono stati uniti i controller: \textit{LoginController}, \textit{UtenteController}, \textit{MenuController} e \textit{TavoloController} in un unico controller; questo perchè non è più necessario suddividere in endpoint multipli le richieste e perché tutti i controller elencati dichiaravano la stessa dipendenza dal servizio \textit{ClientService}.
  \item \textbf{Riscrittura del codice dei resolver}: sono stati riscritte le logiche dei vari resolver seguendo la struttura implementativa precedente; questo perché dopo aver apportato delle modifiche ai tipi in Java è stato necessario di conseguenza modificare alcune logiche;
  \item \textbf{Aggiunta subscription}: essendo state aggiunte due subscription si è reso necessario implementare i resolver corrispondenti.
\end{itemize}
\subsection{Migrazione del frontend da REST a GraphQL}
La migrazione del frontend è stata più semplice e ha seguito le seguenti fasi:
\begin{itemize}
  \item \textbf{Creazione interfacce dei tipi}: sono state create le interfacce per i nuovi tipi strutturati definiti durante la migrazione del backend;
  \item \textbf{Migrazione dei servizi}: è stata effettuata la riscrittura dei servizi in quanto responsabili delle chiamate al GraphQL Server e non più al REST Server;
  \item \textbf{Migrazione dell'interceptor}: l'interceptor è stato riscritto per l'intercettazione delle chiamate GraphQL e non più REST;
\end{itemize}
\subsubsection*{Interfacce dei tipi}
Le interfacce dei tipi sono necessarie per poter creare l'oggetto ricevuto dal backend nel'ambiente di frontend. Sono state dunque aggiunte le seguenti interfacce non presenti precedentemente:
\begin{itemize}
  \item \textit{LoginSession}: precedentemente venivano estrapolati i dati riguardanti l'idToken e la scadenza direttamente dal JSON, ora viene ritornato questo oggetto contenente i medesimi dati;
  \item \textit{OperationSuccess}: precedentemente il codice di stato della risposta HTTP era sufficiente per conoscerne l'esito; questo in GraphQL non accade, dunque viene ritornato appositamente questo oggetto;
  \item \textit{Sessione}: precedentemente veniva direttamente estrapolato dal JSON l'id della sessione, ora viene ritornato questo oggetto dal backend;
\end{itemize}
\subsubsection{Migrazione dei servizi}
I servizi si occupano della comunicazione con il backend per questo motivo sono stati riscritti utilizzando al posto dell'oggetto \textit{HttpClient} per costruire chiamate REST, l'oggetto \textit{Apollo} per costruire quelle GraphQL.
I servizi che sono stati migrati sono:
\begin{itemize}
  \item \textbf{MenuService}: utilizzato per la gestione dei menu;
  \item \textbf{ordiniService}: utilizzato per la gestione degli ordini;
  \item \textbf{authService}: utilizzato per la gestione e verifica dei login;
  \item \textbf{tavoloService}: utilizzato per la gestione delle sessioni nei tavoli;
\end{itemize}
Gli altri servizi presenti non sono stati modificati in quanto non gestivano alcun tipo di comunicazione con il backend.
\subsubsection*{Migrazione dell'interceptor}
L'interceptor di Angular viene creato attraverso l'implementazione dell'interfaccia \textit{HttpInterceptor} ed è stato utilizzato nel frontend dell'applicativo per l'aggiunta nell'header della chiamata del token di identificazione. Il token di identificazione viene ritornato al client dopo il login di un utente e sarà necessario in seguito per identificare le azioni dell'utente loggato.\\ \\
É necessario dunque apportare una modifica all'interceptor realizzato precedentemente poiché non sono più presenti endpoints multipli, ma dopo la migrazione in GraphQL è sufficiente controllare e dunque intercettare le chiamate indirizzate ad un singolo endpoint.
             % Concept Preview
% !TEX encoding = UTF-8
% !TEX TS-program = pdflatex
% !TEX root = ../tesi.tex

%**************************************************************
\chapter{Analisi comparativa dei protocolli REST e GraphQL}
\label{cap:analisi-comparativa}
%**************************************************************
\intro{In questo capitolo verrà svolta l'analisi comparativa tra i due protocolli REST e GraphQL, sia dal punto di vista teorico che da quello pratico .}\\
%**************************************************************
%\section{Analisi comparativa teorica}
%\label{sec:analisi-comparativa-teorica}
%Per analisi comparativa teoria s'intende tutti quegli aspetti che differenziano i due protocolli GraphQL e REST
\section{Introduzione}
REST è stato ed è tutt'oggi lo standard più seguito per la realizzazione delle Web API, tuttavia dopo l'uscita di GraphQL il suo monopolio è stato messo in crisi. Infatti GraphQL ha portato con se delle interessanti soluzioni per molti dei problemi e dei vincoli dello stile architetturale REST.\\
L'innovazione portata da GraphQL è stata apprezzata in larga scala tra gli sviluppatori, a conferma di ciò è possibile visualizzare nell'immagine \ref{graphQL-usage-chart} l'aumento nell'utilizzo di questa tecnologia con il passare degli anni.
\FloatBarrier
\begin{figure}[!ht]
\centering
\includegraphics[width=1\linewidth]{immagini/GraphQLUsageChart.png}
\caption{Grafico sull'utilizzo di GraphQL negli anni.}
\label{graphQL-usage-chart}
\end{figure}
\FloatBarrier
Nel seguente capitolo verranno analizzati nel dettaglio e paragonati i due protocolli, sotto tutti i punti di vista, mettendo in risalto vantaggi e svantaggi di ciascuno; infine verranno riportate le deduzioni elaborate durante lo stage sul protocollo che è meglio seguire in base all'applicativo che si vuole sviluppare.
\section{Analisi comparativa}
Come sottolineato in precedenza, GraphQL e REST hanno diversi aspetti che li differenziano. La più grande differenza tra questi due protocolli è legata alla loro natura: infatti quando si fa riferimento a REST, si sta parlando di uno stile architetturale, dunque di un modo di costruire le proprie API le quali, se rispettano i vincoli REST illustrati al punto \ref{principi-REST}, vengono definite RESTful. Quando si fa riferimento a GraphQL invece, si sta parlando di un linguaggio di query fortemente tipizzato.\\\\
Di seguito è presente un'analisi comparativa dettagliata per ciascun aspetto che differenzia i due protocolli di data fetching.
\subsection{Endpoints}
La prima grossa differenza tra i due protocolli riguarda gli endpoint.\\
Lo stile REST prevede l'utilizzo di più endpoint, sfrutta infatti la molteplicità degli endpoint per differenziare le richieste possibili. Quando un client implementa una richiesta a delle REST API deve sapere esattamente a quale endpoint inviare la richiesta per ricevere i dati necessari. In figura \ref{REST-endpoints} viene rappresentato un REST server che fornisce diversi endpoint e il client, il quale invia diverse richieste a diversi endpoint.
\FloatBarrier
\begin{figure}[!ht]
\centering
\includegraphics[width=1\linewidth]{immagini/RESTEndpoints.pdf}
\caption{Gli endpoints multipli in REST.}
\label{graphQL-usage-chart}
\end{figure}
\FloatBarrier
Con GraphQL questo non avviene, infatti GraphQL prevede l'esposizione di un unico endpoint. A questo singolo endpoint possono essere inviate tutte le richieste inserendo nel body della richiesta la query, la mutation o la subscription. In figura \ref{GraphQL-endpoint} è possibile visualizzare come il GraphQL server fornisca un unico endpoint e come il client invii tutti i tipi di richieste allo stesso medesimo endpoint.
\FloatBarrier
\begin{figure}[!ht]
\centering
\includegraphics[width=1\linewidth]{immagini/GraphQLEndpoint.pdf}
\caption{Il singolo endpoint GraphQL.}
\label{GraphQL-endpoint}
\end{figure}
\FloatBarrier
A proposito di ciò viene riportata di seguito una citazione di Lee Byron, il co-creatore di GraphQL:
\begin{quoting}
  \textit{"Think in graphs, not endpoints."}
\end{quoting}
\subsection{Overfetching e Underfetching}
La questione dell'overfetching e underfetching è uno degli aspetti che viene maggiormente considerato nella decisione architetturale riguardo quale protocollo di data fetching utilizzare tra REST e GraphQL.\\
Lo stile architetturale REST non prevede di definire lato client esattamente quali dati ricevere. Un client che necessita un certo insieme da un server con REST API, è costretto ad eseguire una o più richieste e dunque a prendersi in carico della rielaborazione dei dati. Per uno sviluppatore backend è molto complesso riuscire a creare delle REST API che siano in grado di soddisfare esattamente le tutte le richieste dei client.
L'introduzione di GraphQL come nuovo protocollo di data fetching ha posto una soluzione a questo problema, permettendo al client di specificare esattamente la forma e il quantitativo di dati necessari. Quando si decide di implementare un applicativo e si valuta quale protocollo di data fetching utilizzare, questo è sicuramente un punto da considerare.\\
\paragraph{Overfetching}
Quando si parla di overfetching si fa riferimento al fatto che vengano forniti più dati di quanti realmente necessari. Riprendendo come esempio il prototipo visto nel capitolo \ref{casi-uso}, si suppone che il client necessiti della lista degli impiegati e che, per ciascun impiegato, necessiti  esclusivamente l'id e il nome. Qualora si tratti della version REST del server, il client inviare la richiesta HTTP all'endpoint mappato sul metodo \textit{allEmployee()}, il quale ritorna una lista di impiegati e, per ciascun impiegato, tutti i campi. Di seguito il JSON di risposta che il client riceve in seguito alla richiesta nel caso in cui fossero presenti solo due impiegati:
\begin{verbatim}
  [
    {
        "id": "3AFASDF12F",
        "name": "Matteo",
        "surname": "Verdi",
        "salary": 1500,
        "birth": "1995-02-21"
    },
    {
        "id": "GA14PL3FAV",
        "name": "Marco",
        "surname": "Blu",
        "salary": 1500,
        "birth": "1993-12-20"
    }
  ]
\end{verbatim}
Si può subito notare come i campi \textit{surname}, \textit{salary} e \textit{birth} non siano necessari in quanto il client utilizza solo i campi \textit{id} e \textit{name}. Nel caso del prototipo si tratta di un problema irrisorio data la ridotta grandezza dei dati, tuttavia in applicativi a larga scala che richiedono grossi quantitativi di dati complessi può risultare un problema in termini di occupazione di rete e di rallentamenti dell'applicativo. Questo problema può essere risolto in un server con REST API fornendo endpoint specifici per ciascun tipo di richiesta, tuttavia questa soluzione rischia di introdurre confusione e nel tempo la manutenzione potrebbe risultare sempre più complessa.\\
GraphQL risolve questo problema attribuendo al client la responsabilità di definire quali siano i campi di cui necessita. Questo è possibile specificando nella query i campi richiesti, di seguito l'esempio dell'invocazione contenente a sinistra la query \textit{allEmployee}, mentre a destra il JSON ricevuto in risposta:
\begin{verbatim}
    QUERY                                     JSON RETURNED

    query getAllEmployee{                 [
      allEmployee {                         {
        Employee {                              "id": "3AFASDF12F",
          id                                    "name": "Matteo"
          name                              },
          }                                 {
      }                                          "id": "GA14PL3FAV",
    }                                            "name": "Marco"
                                            }
                                          ]
\end{verbatim}
In GraphQL questo approcio permette di non sovraccaricare inutilmente la rete e di mantenere ordinato e di facile manutenzione la struttura API del server. Lato client è richiesto un maggior sforzo per nella specifica della query, ma d'altra parte si evitano problemi di rallentamento o errori dovuti al fetching di grossi quantitativi di dati inutili. 
\section{Conclusioni}





% Spiegazione di cosa si intende per analisi comparativa teorica: analisi % comparativa che va a comparare gli aspetti prettamente teorici, come ad esempio:
% \begin{itemize}
%  \item overfetching e underfetching;
%  \item un protocollo è una sorta di standard/stile architetturale (restfull api)     mentre l'altro un linguaggio di query fortemente tipizzato (REST invece non è safe dal punto di vista dei tipi);
%  \item come sfruttano il protocollo http (stati risposte http, endpoint multipli vs singolo, ecc...);
%  \item manuntenibilità nel tempo;
%  \item documentazione (GraphQL si autodocumenta, REST no);
%  \item meccanismo di caching integrato (mancante in GraphQL);
%  \item formati output di risposta (GraphQL --> JSON, REST --> JSON, XML, YAML);
% \end{itemize}
% \section{Analisi comparativa sui casi d'uso}
% Analisi comparativa basata su:
% \begin{itemize}
%  \item differenze nell'analisi e progettazione iniziale delle API;
%  \item differenze durante lo sviluppo delle API dal punto di vista di BE e FE(anche legate agli strumenti utilizzati ad es. Spring Data REST vs Spring GraphQL);
%  \item differenze prestazionali (utilizzato tool K6 per load test);
% \end{itemize}2
             % Product Prototype
% !TEX encoding = UTF-8
% !TEX TS-program = pdflatex
% !TEX root = ../tesi.tex

%**************************************************************
\chapter{Tecnologie utilizzate}
\label{cap:tecnologie}
%**************************************************************
             % Product Design Freeze e SOP
\input{capitoli/capitolo-7}             % Conclusioni
% \input{capitoli/capitolo-8}
\appendix
\input{capitoli/capitolo-A}             % Appendice A

%**************************************************************
% Materiale finale
%**************************************************************
\backmatter
\printglossary[type=\acronymtype, title=Acronimi e abbreviazioni, toctitle=Acronimi e abbreviazioni]
\printglossary[type=main, title=Glossario, toctitle=Glossario]
\input{inizio-fine/bibliografia}
\end{document}
